% concatenated from arxiv.tex
\documentclass{amsart}
\usepackage{xcolor,enumerate}       % blackboard math symbols
\topmargin  = 0.0 in
\leftmargin = 0.9 in
\rightmargin = 1.0 in
\evensidemargin = -0.10 in
\oddsidemargin =  0.10 in
\textheight = 8.5 in
\textwidth  = 6.6 in
\setlength{\parskip}{2mm}
\setlength{\parindent}{0mm}
\def\baselinestretch{1}
 \definecolor{darkgreen}{rgb}{0,0.5,0}
\usepackage{blindtext}
\usepackage{mathtools}
\usepackage{tikz}
\usetikzlibrary{arrows.meta,calc,positioning,decorations.pathmorphing}

%\usepackage[utf8]{inputenc} % allow utf-8 input
\usepackage[T1]{fontenc}    % use 8-bit T1 fonts
\usepackage{hyperref}       % hyperlinks
\usepackage{url}            % simple URL typesetting
\usepackage{amsmath,amssymb}
\usepackage{amsmath,mathrsfs,bm}
\usepackage{parskip}
\usepackage{paralist}
\usepackage{hyperref}
\usepackage{bigints}
\usepackage{graphicx}     % For including images
\usepackage{subcaption}   % For subfigures with individual captions\textbf{}
\usepackage{caption}
\usepackage{graphicx} % Required for inserting images
\usepackage{comment}
\usepackage{mathrsfs}

\usepackage{lipsum}
\usepackage{amsfonts}
\usepackage{epstopdf}
%\usepackage{algorithmic}
%\usepackage{graphicx}     % For including images
\ifpdf
  \DeclareGraphicsExtensions{.eps,.pdf,.png,.jpg}
\else
  \DeclareGraphicsExtensions{.eps}
\fi

\usepackage{algorithm}       % Floating environment
\usepackage{algpseudocode}   % Pseudocode formatting

\newtheorem{theorem}{Theorem}[section]
\newtheorem{conjecture}[theorem]{Conjecture}
\newtheorem{lemma}[theorem]{Lemma}
\newtheorem{claim}[theorem]{Claim}
\newtheorem{problem}[theorem]{Problem}
%\newtheorem{algorithm}[theorem]{Algorithm}
\newtheorem{proposition}[theorem]{Proposition}
\newtheorem{corollary}[theorem]{Corollary}
\newtheorem{construction}[theorem]{Construction}
\newtheorem{definition}[theorem]{Definition}
\newtheorem{example}[theorem]{Example}
\newtheorem{xca}[theorem]{Exercise}
\newtheorem{comments}[theorem]{Comments}
\newtheorem{remark}[theorem]{Remark}
\newtheorem{addendum}{Addendum}
\newtheorem{intuition}{intuition}
\newtheorem{assumption}{Assumption}
%\newtheorem{example}{Example}
%\newtheorem{problem}{Problem}

\numberwithin{equation}{section}

\newcommand{\R}{\mathbb{R}}
\newcommand{\map}[3]{#1:#2 \rightarrow #3}
\newcommand{\E}{\mathbb{E}}
\newcommand{\KL}{D_{\mathrm{KL}}}
\newcommand{\PP}{\mathrm{P}}
\newcommand{\PPu}{\PP^u}
%\newcommand{\E}{\mathbb{E}}
\newcommand{\argmin}{\mathop{\mathrm{arg min}}}
\newcommand{\argmax}{\mathop{\mathrm{arg max}}}
\newcommand{\kecmt}[1]{{\color{blue} #1}}
\newcommand{\SE}{\mathrm{SE}}
\newcommand{\diff}{\mathrm{d}}
\newcommand{\e}{\mathrm{e}}
\newcommand{\ii}{\mathrm{i}}
\newcommand{\tr}{\mathrm{tr}}
\newcommand{\Z}{\mathbb{Z}}
\newcommand{\calV}{\mathcal{V}}
\newcommand{\calF}{\mathcal{F}}
\newcommand{\calL}{\mathcal{L}}
\newcommand{\inner}[2]{\langle #1 , #2 \rangle}
\newcommand{\LL}{\mathcal{L}}




\begin{document}

\begin{abstract}
We study the least-energy way to reshape a probability distribution when motion is constrained to a horizontal bundle, i.e., optimal transport and distribution steering in sub-Riemannian geometry—motivated by density control over underactuated systems. To obtain a continuous and numerically tractable formulation, we introduce an entropic regularization by adding small noise aligned with the control directions and study the associated Schr\"odinger bridge problem. The resulting reference process is a degenerate diffusion on the sub-Riemannian manifold. Under bracket-generating hypotheses we obtain smooth, strictly positive transition densities and a forward--backward characterization of the optimal bridge. This leads to a practical Sinkhorn-type algorithm for the Schr\"odinger potentials and, as the noise level vanishes, a recovery of the deterministic sub-Riemannian optimal transport problem. We demonstrate with a numerical example.
\end{abstract}

\title[From Schr\"odinger Bridge to Optimal Transport over Sub-Riemannian Manifolds]{From Schr\"odinger Bridge to Optimal Transport over Sub-Riemannian Manifolds}

%\thanks{The work of the author is supported by the Natural Sciences and Engineering Research Council of Canada.} 


\author[D. O. Adu, K. Elamvazhuthi, and B. Gharesifard ]{D. O. Adu, K. Elamvazhuthi, and B. Gharesifard}
%\address{Department of Mathematics and Statistics\\
%Queen's University\\
%Kingston, ON, Canada}
%\email{bahman.gharesifard@queensu.ca}
%\urladdr{https://gharesifard.github.io}

\maketitle

\section{Introduction}\label{sec:intro}
Motivated by underactuated vehicles~\cite{murray2017mathematical,lee2010geometric}, minimum effort control over Sub-Riemannan manifolds~\cite{tiwariJena2020optimal,tiwari2022optimal,halder2021sub} and the need to reshape an entire distribution rather than single trajectories~\cite{brockett2007optimal,agrachev2009optimal,chen2021optimal,adu2024schrodinger,hindawi2011mass}, we ask the following question: what is the minimum-effort feedback law that steers the ensemble of trajectories of a control-affine system from an initial to a prescribed terminal distribution? Usually, initial data comprises an initial and a terminal probability measure $\mu_0,\mu_f\in\mathcal{P}(\mathbb{R}^d)$ with densities $\rho_0,\rho_f$ on $\mathbb{R}^d$ and a time-varying matrix-valued map $\map{g}{[0,t_f]\times\mathbb{R}^d}{\mathbb{R}^{d\times m}}$. The goal is to find a control $u$, if it exists, that transports with minimum effort along the control-affine system
\begin{equation}\label{eq:intro_det_dyn}
  \dot{x}_t^u  =  g(t, x_t^u)u(t, x_t^u), \qquad  x_0^u \sim \mu_0,\quad t\in[0,t_f].
\end{equation}
Here, $x^u_t\in\mathbb{R}^{d}$ denotes the state process influenced by the time-varying feedback control $\map{u}{[0,t_f]\times\mathbb{R}^m}{\mathbb{R}^{d}}$. The notation $x_0^u \sim \mu_0$ implies that at $t=0$ the random variable $x_0^u$ is distribution according to $\mu_0$. We say that the system is underactuated whenever $m<d$. We measure effort via %an $\alpha$-power control cost, with $\alpha\in\{1,2\}$ in applications that we consider:
\begin{equation}\label{eq:intro_det_obj}
  \inf_{u\in\mathcal U}\ \E \left[\int_{0}^{t_f}\frac12 \|u(t, x_t^u)\|^2 dt\right]
  \quad\text{s.t.}\quad  x_{t_f}^u\sim\mu_f,
\end{equation}
%\margin{Make this first footnote part of the text. Not good to have a footnote on page one.}
where the expectation is with respect to the law of $x^u$ from $ x_0^u\sim\mu_0$ to $x_{t_f}^u\sim\mu_f$. Here a convenient admissible class is 
%\margin{B: make the way you put expectation behind the integral consistent throughout. I fixed this one}
$$\mathcal U:=\Big\{u:\ [0,t_f]\times\mathbb{R}^d\to\mathbb{R}^m:
\E \left[\int_0^{t_f}\|u(t, x_t^u)\|^2 dt\right]<\infty\Big\}.$$
Throughout, we assume that the densities of $\rho_0,\rho_f$ are positive and compactly supported. The latter assumption can be replaced with having an exponential tail on $\R^d$. The structural constraint $m<d$ implies that the admissible velocities $\dot{x}_t$ in~\eqref{eq:intro_det_dyn} are constrained to lie in the distribution $\dot{x}_t\in D_{(t,x_t)}$, where
\begin{equation}\label{eq: time_varying_distrubtion}
\mathcal D_{(t,x_t)}:=\mathrm{span}\{g_1(t,x_t),\dots,g_m(t,x_t)\}\subset T_{x_t}\R^d\equiv\R^d,    
\end{equation}
 with $g_i(t,x_t)$ the $i$-th column of $g(t,x_t)$ in~\eqref{eq:intro_det_dyn}. %\margin{So I like to make sure that: whenever you refer to $g_i$ as a column, you write $g_i(t,x_t)$, as otherwise, this is a vector field}
 This distribution induces a sub-Riemannian structure with a metric $g_x^{\text{SR}}(v,v) := v^\top G(x)^\dagger v$ for $v \in \mathcal{D}_{(t,x)}$, where $G^\dagger$ is the pseudoinverse of $G := g g^\top$. The associated sub-Riemannian distance is  
 \[
 d_{SR}(x_0, x_f) = \inf\left\{\int_0^{t_f} \sqrt{g^{SR}_{x_t}(\dot{x}_t, \dot{x}_t)} dt : \dot{x}_t \in \mathcal{D}_{(t,x_t)},~  ~  x_{t_f}=x_f\right\}
 \]
 and using~\eqref{eq:intro_det_obj}, we have that 
\begin{equation*}
%c_2(x_0, x_f) = 
\frac12d^2_{SR}(x_0, x_f)= \inf\left\{\int_0^{t_f} \frac12\|\dot{x}_t\|_{G^\dagger}^2 dt : \dot{x}_t \in \mathcal{D}_{(t,x_t)},~ ~  x_{t_f}=x_f\right\},
\end{equation*}
where $\|v\|_{G^\dagger}^2 = v^\top G^\dagger v$. Therefore, the optimal affine control~\eqref{eq:intro_det_dyn}-\eqref{eq:intro_det_obj} is a dynamic ensemble control problem over a sub-Riemannian manifold $(\R^d,\mathcal{D},g^{SR})$, which seeks a feedback control law $u(t,x)$ steering $\mu_0$ to $\mu_f$ and the dynamics of the ensemble are constrained to the distribution $\mathcal{D}$ in~\eqref{eq: time_varying_distrubtion}. This problem has been studied in~\cite{elamvazhuthi2023dynamical,elamvazhuthi2024benamou}, where the authors established the existence of optimal feedback control using Young measures for general nonlinear control affine-systems~\cite{elamvazhuthi2023dynamical} and subsequently, as regular feedback control function under suitable absence-of-abnormality assumptions~\cite{elamvazhuthi2024benamou} leveraging results of \cite{agrachev2009optimal,figalli2010mass}. At times, we refer to problem~\eqref{eq:intro_det_dyn}-\eqref{eq:intro_det_obj} as an optimal transport problem over a sub-Riemannian manifold, because it induces a static optimal transport (OT) problem over a sub-Riemannian manifold
\begin{equation}\label{eq:OT-static}
 \inf_{\gamma\in\Pi(\mu_0,\mu_f)} \int_{\mathbb{R}^d\times\mathbb{R}^d} \frac12d_{SR}^2(x_0, x_f) , d\gamma(x_0,x_f),
\end{equation}
where
\begin{equation}\label{eq: transport_plans}
\Pi(\mu_0,\mu_f):=\{\gamma\in\mathcal P(\R^d\times\R^d) : \pi_{0\#}\gamma =\mu_0\quad\text{and}\quad \pi_{f\#}\gamma =\mu_f\}   
\end{equation}
is the set of couplings of $(\mu_0,\mu_f)$ and $\pi_0(x_0,x_f)=x_0$ and $\pi_f(x_0,x_f)=x_f$, for all $x_0,x_f\in\R^d$, are the projection maps. This static problem ~\eqref{eq:OT-static}-\eqref{eq: transport_plans}, which seeks the minimal cost coupling between $\mu_0$ and $\mu_f$, has been studied in~\cite{agrachev2009optimal,figalli2010mass}. %  for general nonholonomic systems and by~\cite{figalli2010mass}  for sub-Riemannian manifolds.} 
\begin{comment}
formulated in~\cite{agrachev2009optimal}, where for a fixed pair $(x_0,x_f)$, the transport cost is set as 
\begin{equation}\label{eq:endpoint-cost}
  c(x_0,x_f)
  :=
  \inf\Big\{
    \int_0^{t_f}\frac12 \|u(t)\|^2 dt
  \Big\},
\end{equation}
{\color{blue}subject to~\eqref{eq:intro_det_dyn}, where $x^u_0=x_0$ and $x^u_{t_f}=x_f$ almost surely,
with $c(x_0,x_f)=+\infty$ if no admissible pair exists (lack of
reachability)}.
% \begin{comment}
% {\color{red} which measures the system
% \[ 
%     \dot x_t=g(t,x_t)u(t),\qquad x(0)=x_0\quad \text{and}\quad x(t_f)=x_f
% \]
% with the convention that $c(x_0,x_f)=+\infty$ if no admissible pair exists (lack of
% reachability).}\margin{What do we mean by measures? Also, we should simply refer to the governing equations we had before. This is written differently} 
% \end{comment}
As usual
 The control system induces a sub-Riemannian structure on $\mathbb{R}^d$
with horizontal distribution $\mathrm{span}\{g_1,\dots,g_m\}$, where $g_i$ is the $i$-th column of $g$ in~\eqref{eq:intro_det_dyn}. The cost \eqref{eq:intro_det_obj} with $\alpha=2$ corresponds to minimizing
the {\color{blue}energy} %{\color{red} sub-Riemannian (SR) energy}\margin{Not a common terminology, I think, and we have not said anything about sub-Riemannian yet. I would just say energy}:
\[
E[\gamma] = \frac12\int_0^{t_f} \|\dot{\gamma}(t)\|_{G^{-1}}^2 dt,
\]
%\margin{Why would $ G $ be invertible? We also need to define the norm we use here (standard, but still)}
where {\color{blue}$\dot{\gamma}\in\mathrm{span}\{g_1,\dots,g_m\}$, }.  %\margin{You need to use $ G^\dagger $ then! }
In this case, for $\gamma(0)=x_0$ and $\gamma(t_f)=x_f$, we get that
\begin{equation}\label{eq: Riemannian_distance_transport_cost}
c_2(x_0,x_f)=\tfrac12 d_g(x_0,x_f)^2,    
\end{equation}
{\color{blue}where $d_g$ denotes the sub-Riemannian distance induced by the distribution}
see, e.g., \cite[Chs.~2–3]{Montgomery2002}, \cite[Chs.~3,~15]{AgrachevSachkov2004}, and \cite{Bellaiche1996,Jurdjevic1997}. 
For $\alpha=1$ one obtains a \emph{sub-Finsler} (``minimum fuel'') distance associated with a norm on the {\color{blue} horizontal distribution}, see \cite[\S15.6]{AgrachevSachkov2004} and the survey \cite{LeDonne2017}. For 

The static OT problems,~\cite{agrachev2009optimal,rifford2012sub,figalli2010mass} have established existence of Monge solutions for nonlinear control problems and \cite{hindawi2011mass} considered the linear quadratic case. The dynamic OT setup in~\eqref{eq:intro_det_dyn}--\eqref{eq: Riemannian_distance_transport_cost} 
\end{comment}
As stated in the literature~\cite{agrachev2009optimal,figalli2010mass,elamvazhuthi2023dynamical,elamvazhuthi2024benamou,rifford2012sub}, the core difficulty of OT over a sub-Riemannian manifold stems from the non-smooth geometry induced by the constraints: the control-to-state map can have singularities, so the regularity and uniqueness of minimizers may fail; even when a Monge map exists, it can be discontinuous or non-unique. The novelty in this work is to establish a framework that mitigates these difficulties.

Motivated by the regularizing effects of entropic penalization in optimal transport \cite{cuturi2013sinkhorn}, our contribution is of a different nature: we introduce an \emph{entropic regularization} of sub-Riemannian optimal transport. Specifically, we perturb~\eqref{eq:intro_det_dyn} by a small diffusion acting \emph{only along the horizontal distribution}:
\begin{equation}\label{eq:intro_sto_dyn}
dX_t^u = g(t,X_t^u) u(t,X_t^u) dt + \sqrt{\epsilon} g(t,X_t^u) dW_t,\qquad X_0^u\sim\mu_0,
\end{equation}
where $W_t$ is the standard $\R^m$-Brownian motion, with $\epsilon>0$, and consider the problem
\begin{equation}\label{eq:intro_sto_obj}
\inf_{u\in\mathcal U}\ \E\!\left[\int_0^{t_f}\frac12\|u(t,X_t^u)\|^2dt\right]
\quad\text{subject to}\quad X_{t_f}^u\sim\mu_f.
\end{equation}
This stochastic control problem is equivalent to the Schr\"odinger bridge problem (SBP)~\cite{schrodinger1931umkehrung}: among all path measures on $\Omega:=C([0,t_f];\R^d)$ that are absolutely continuous with respect to a reference law and match the endpoint marginals $(\mu_0,\mu_f)$, find the one that minimizes the relative entropy with respect the reference law. A difficulty immediately arises: the reference law, if it exists, is degenerate and the setup~\eqref{eq:intro_sto_dyn}-\eqref{eq:intro_sto_obj} is ill-posed. The associated Kolmogorov operator is not uniformly elliptic, and parabolic regularity theory in~\cite{dai1991stochastic,Leonard2012} does not apply directly.  To define a well-defined   Schr\"odinger system, one requires a strictly positive transition density function. Unlike the linear stochastic degenerate case~\cite{adu2022optimal,adu2024schrodinger,chen2021optimal,chen2016optimal} where the presence of drift is helpful, for a non-linear case, establishing the existence of a positive smooth transition density is delicate. While uniformly ellipticity ensures smooth and everywhere positive transition density function, these desirable properties are decoupled in the non-linear non-elliptic diffusion case. For smoothness, one typically imposes the H\"ormander condition which leads to hypoellipticity of the generator~\cite{Hormander1967}. Under H\"ormander, positivity is tied to submersion property of the endpoint map~\cite{leandre2005positivity,arous1991decroissance} (Bismut/submersion condition) rather than the Chow-Rashevskii controllability condition. The driftless system is the structural choice that ensures that Chow-Rashevskii controllability is equivalent to H\"ormander-Bismut/submersion condition~\cite{arous1991decroissance}. Aside the fact that our settings smoothens the geometric singularities present in the OT-SR settings, the SBP setup is strictly convex and hence acts as a  selection principle up a to subsequence among potentially many deterministic SR-optimal solutions. Moreover, the dynamic regularity formulation transforms the continuity equation associated to~\eqref{eq:intro_det_dyn} 
to a parabolic  associated to~\eqref{eq:intro_sto_dyn}. This adds a numerical novelty in the paradigm of OT-SR setup. Aside the fact that the hypoelliptic forward-backward characterization improves 
numerical stability, it yields a practical algorithmic pipeline that mirrors (and extends) the classical entropic optimal transport/Sinkhorn 
paradigm to the sub-Riemannian regime. We show that the Schr\"odinger system associated to our problem can be solved via a Sinkhorn/IPFP 
fixed-point iteration on the cone of positive functions. Finally, the hypoelliptic degeneracy introduces a genuine 
discretization challenge absent in uniformly elliptic settings: standard Euclidean 
finite-difference discretizations may fail to preserve positivity and mass. Our 
schemes therefore, discretize the generator and its adjoint in a manner that is consistent 
with the underlying sub-Riemannian geometry. 


We state here that L\'eonard in~\cite{Leonard2012} studied the Schr\"odinger problem in a general Markovian-process framework and in particular, analyzed the diffusion-to-optimal-transport limit in the classical Brownian/elliptic setting. By contrast, our work treats a genuinely degenerate Markovian process, where the noise is \emph{anisotropic} and the zero-limit yields sub-Riemannian optimal transport.
%{\color{blue} We state here that convergence of entropic Schr\"odinger bridges to the deterministic Monge--Kantorovich problem has been studied systematically by L\'eonard in~\cite{Leonard2012} under the assumption that the reference diffusion is uniformly elliptic.} By contrast, our work introduces an \emph{anisotropic} entropic regularization: the noise acts only along the horizontal distribution, yielding a degenerate reference process whose law is supported on a null set for the full Wiener measure in~\cite{Leonard2012}. This horizontal regularization preserves the sub-Riemannian geometry and leads, in the zero-noise limit, to the sub-Riemannian optimal transport problem. %, removes the ellipticity hypothesis and establish $\Gamma$-convergence to the zero-noise limit on the path space.
%\lipsum[2-3]

% The outline is not required, but we show an example here.
The paper is organized as follows: In Section~\ref{sec: Stochastic Geometric Control Problem}, we discuss the well-posedness of the stochastic optimal control problem via Schr\"odinger bridge problem. In Section~\ref{sec:Eulerian} we derive the first-order condition of optimality in the Eulerian framework. In Section~\ref{sec: Limiting Section}, we show that the stochastic solution established in Section~\ref{sec:Eulerian} converges to the corresponding Eulerian solution in Section~\ref{sec: Stochastic Geometric Control Problem}. Section~\ref{sec: Numerical Method} is dedicated to numerical method and example.

\section{Stochastic Geometric Control Problem}\label{sec: Stochastic Geometric Control Problem}
We consider the problem of transporting an initial particle distribution $\mu_0$ to a final $\mu_f$ via a controlled dynamical system. The velocities of the ensemble of particles initially distributed according to $\mu_0$ are restricted to a prescribed subspace at every point. The objective is to achieve this transport while minimizing the total expected control effort, leading to an optimal transport problem with geometric constraints.

%\margin{$\mathcal{B}$ undefined} 
We now formalize this problem. Our aim is to investigate the optimal distributional control problem:
\begin{equation}\label{eq: expectation_total_control_effort}
\inf_{u} \E \left[ \int_0^{t_f} \frac12 \|u(t,  x_t)\|_2^2  dt\right]
\end{equation}
subject to
\begin{equation}\label{eq: Num_boundary condition}
dX_t=g(x_t)u(t,x_t)dt,\quad x_0\sim\mu_0\quad\text{and}\quad x_{t_f}\sim\mu_f.
\end{equation}
To simplify the notation, from this point on, we suppress the explicit $t$-dependence and write $g(x)$; the analysis remains identical for $g(t,x)$ as in~\eqref{eq:intro_det_dyn}. Thus, the problem is to find the minimum expected energy required to transport from  $\mu_0$  to  $\mu_f$, where $\mu_0$ and $\mu_f$ have compactly supported strictly positive densities $\rho_0$ and $\rho_f$, respectively, along admissible paths whose velocities lie in the distribution%~\eqref{eq: distribution}
\begin{equation}\label{eq: distribution}
\mathcal{D} := \operatorname{span}\{g_1, \dots, g_m\} \subset T\mathbb{R}^d.
\end{equation}
We emphasize that, the $\mathcal{D}$ above equipped with the metric  $g_x^{\text{SR}}(v,v) := v^\top G(x)^\dagger v$ for $v \in \mathcal{D}_{x}$, where $G^\dagger$ is the pseudoinverse of $G := g g^\top$ constitutes the sub-Riemannian structure $(\R^d,\mathcal{D},g^{SR})$ on the state space. Accoringly, we refer to~\eqref{eq: expectation_total_control_effort}-\eqref{eq: Num_boundary condition} optimal transport over sub-Riemannian manifold $(\R^d,\mathcal{D},g^{SR})$. Henceforth we denote this geometric structure simply by $(\R^d,\mathcal{D},g^{SR})$.

To ensure the problem is well-posed, we make the following two assumptions.
\begin{assumption}\label{ass:wellposed}
The coefficients $g_i \in C^\infty(\mathbb{R}^d; \mathbb{R}^d)$, where $i\in\{1,\dots.m\}$, has bounded derivatives of all order. 
\end{assumption}
The second is the following.
\begin{assumption}\label{ass:Controllability condition}
The vector fields ${g_1, \dots, g_m}$ satisfy the bracket-generating condition. Specifically, we assume that the Lie algebra generated by $\{g_1, \dots, g_m\}$ spans the tangent space $T_y\mathbb{R}^d$ at every point $y \in \mathbb{R}^d$:
\[
\operatorname{Lie}(g_1, \dots, g_m)(y) = \mathbb{R}^d \quad \text{ for all $y \in \mathbb{R}^d$.}
\]
\end{assumption}

We emphasize here that, all the analysis in this work is under the umbrella of Assumption~\ref{ass:wellposed} and~\ref{ass:Controllability condition}. 
Rather than attempting to solve~\eqref{eq: expectation_total_control_effort}-\eqref{eq: Num_boundary condition} directly as in \cite{agrachev2009optimal,figalli2010mass,elamvazhuthi2024benamou}, we first attempt to solve the relaxed problem: %stochastic optimal control problem
\begin{equation}\label{eq: objective}
\inf_{u \in \mathcal{U}} \E \left[ \int_0^{t_f} \frac12 \|u(t,  X_t^u)\|_2^2  dt  \right]
\end{equation}
subject to 
\begin{align}\label{controlled_dynamics}
&d X_t^u = (b_{\epsilon}(x_t^u) + g(X_t^u) u(t,  X_t^u))  dt+ \sqrt{\epsilon} g(X_t^u)  dW_t,\\
& X_0^u\sim\mu_0\quad\text{and}\quad  X_{t_f}^u\sim\mu_{t_f},\nonumber
\end{align}
where $b_{\epsilon} :=\frac\epsilon2\sum_{i=1}^{m}\nabla_{g_i}g_i$ %, the derivative of the vector field $ g_i$ in the direction of $ g_i$, 
is the Stratonovich correction drift ensuring that %the uncontrolled SDE 
\begin{equation}\label{uncontrolled_dynamic_constraint}
dX_t = b_{\epsilon}(X_t) dt + \sqrt\epsilon g(X_t) dW_t, \quad  X_0 \sim \mu_0
\end{equation}
naturally stays in the distribution~\eqref{eq: distribution}. %{. See Figure~\ref{fig:diffusion_comparison_three} for a demonstration.}\margin{Not sure this is the place to refer to this?} 
Here, $\{W_t\}_{t \ge 0}$ is the standard $\mathbb{R}^m$ Brownian motion,  
$\epsilon>0$ is the noise intensity, and
\[
\mathcal{U} := \left\{ \map{u}{[0, t_f] \times \mathbb{R}^d}{\mathbb{R}^m} \middle| \text{$u$ is adapted and } \mathbb{E} \int_0^{t_f} \|u(t,  X_t^u)\|^2 dt < \infty \right\}.
\]
Later we show that the solution for~\eqref{eq: expectation_total_control_effort}-\eqref{eq: Num_boundary condition} is obtained as a limiting process of the relaxed problem. Problem~\eqref{eq: objective}-\eqref{controlled_dynamics} is a Schr\"odinger bridge problem over a sub-Riemannian manifold $(\R^d,\mathcal{D},g^{SR})$. To make this document self-contained and also emphasize the novelty of our work in that paradigm, we state the derivation below.
\begin{figure}[t]
    \centering
    
    % Left plot: Horizontally Constrained
    \begin{subfigure}[b]{0.32\textwidth}
        \centering
        \includegraphics[width=\textwidth]{Horizontally_Constrained_Process.png}
        \caption{Horizontally constrained diffusion}
        \label{fig:constrained}
    \end{subfigure}
    \hfill
    % Middle plot: Purely Horizontal
    \begin{subfigure}[b]{0.32\textwidth}
        \centering
        \includegraphics[width=\textwidth]{Pure_Horizontal_Constraint.png}
        \caption{Purely horizontal diffusion}
        \label{fig:horizontal_only}
    \end{subfigure}
    \hfill
    % Right plot: Unconstrained Diffusion
    \begin{subfigure}[b]{0.32\textwidth}
        \centering
        \includegraphics[width=\textwidth]{Unconstrained_Process.png}
        \caption{Isotropic diffusion}
        \label{fig:unconstrained}
    \end{subfigure}
    
    \caption{Plots~(a) and~(b) are anisotropic diffusion $dX_t=\sqrt{\epsilon} g(X_t) dW_t$ in $\mathbb{R}^3$ with $\epsilon=1$ as opposed to~(c) which is an isotropic diffusion $dX_t=\sqrt{\epsilon}dW_t$ with $\epsilon=1$ in $\mathbb{R}^3$. In all scenarios, the same $10$ initial samples from $X_0\sim\mu_0$ are used. The plot in~(a) is a horizontally constrained (hypoelliptic) diffusion of Heisenberg type, encoded by $g$ in Example~\ref{ex: unified_example}, where the noise acts only along the horizontal distribution and the vertical coordinate is generated indirectly through the noncommutative geometry. This creates a relatively flat pattern in the noise. The plot in~(b) corresponds to a process which is purely confined to horizontal planes (encoded by $g$).}
    \label{fig:diffusion_comparison_three}
\end{figure}

\subsection{Schr\"odinger Bridge Formulation}
Let $\Omega:=C([0,t_f];\mathbb R^d)$ be the canonical space  with the uniform topology (and endowed with the Borel $\sigma$-algebra) of the uncontrolled process $\{ X_t\}_{0\leq t}$ in~\eqref{uncontrolled_dynamic_constraint}. Note that under Assumption~\ref{ass:wellposed}, we have that~\eqref{uncontrolled_dynamic_constraint} admits a unique strong solution on $[0,t_f]$ (see e.g.,~\cite{oksendal2003stochastic}). Let $\mathrm{R}^{\epsilon,\mu_0}\in\mathcal P(\Omega)$, satisfying the pushforward \footnote{Given a measurable map $\map{T}{ (\Omega,\mathcal{F})}{(Y,\mathcal{G})}$ 
and a measure $\mu$ on $(\Omega,\mathcal{F})$, 
the \emph{pushforward measure} $T_\# \mu$ on $(Y,\mathcal{G})$ is defined by 
$T_\# \mu(B) := \mu(T^{-1}(B))$ for every $B \in \mathcal{G}$.} 
\begin{equation}\label{eq: law of unreflected process}
 ( \mathrm{ev}_0)_\# \mathrm{R}^{\epsilon,\mu_0} = \mu_0,   
\end{equation} 
 be the canonical law of~\eqref{uncontrolled_dynamic_constraint} on the sub-Riemannian manifold $(\R^d,\mathcal{D},g^{SR})$. Here $\map{\mathrm{ev}_t}{\Omega}{\R^d}$ is the evaluation $\mathrm{ev}_t(X)=X_t$, for all $X\in\Omega$ at time $t\geq 0$. For a given admissible control $u \in \mathcal{U}$, let $\mathrm{P}^u \in \mathcal{P}(\Omega)$ be the corresponding probability distribution of the process $\{ X_t^u\}_{0\leq t\leq t_f}\in \mathbb{R}^d$ in~\eqref{controlled_dynamics}. The Schr\"odinger bridge problem (SBP) is the infinite-dimensional optimization problem: 
\begin{equation}\label{sbp_precise}
\inf_{\mathrm{P}^u \in \mathcal C_{\mathrm{SB}}} \KL(\mathrm{P}^u \| \mathrm{R}^{\epsilon,\mu_0}):= \mathbb{E}_{\mathrm{P}^u} \Bigg[ \log \frac{d\mathrm{P}^u}{d\mathrm{R}^{\epsilon,\mu_0}} \Bigg] 
\end{equation}
where
\[
\mathcal C_{\mathrm{SB}}
:= \Big\{\mathrm{P}\ll\mathrm{R}^{\epsilon,\mu_0} :\ 
(\mathrm{ev}_0)_\#\mathrm{P}=\mu_0,\ (\mathrm{ev}_{t_f})_\#\mathrm{P}=\mu_f\Big\}.
\]

By the Girsanov transformation~\cite[Chapter~8.6]{oksendal2003stochastic},
\begin{equation}\label{rn_derivative_precise}
\frac{d\mathrm{P}^u}{d\mathrm{R}^{\epsilon,\mu_0}}\Big|_{\mathcal{F}_t} 
= \exp\Bigg( \int_0^t \frac{1}{\sqrt{\epsilon}} u(s,  X_s)^\top dW_s - \frac12 \int_0^t \frac{1}{\epsilon} \|u(s,  X_s)\|^2 ds \Bigg). 
\end{equation}
Therefore, under $\mathrm{P}^u$, the distribution of \eqref{controlled_dynamics},
\[
W_t^u := W_t - \int_0^t \frac{1}{\sqrt{\epsilon}} u(s,  X_s) ds
\]
is a Brownian motion. 
Finally, using~\eqref{rn_derivative_precise}, the objective functional in~\eqref{sbp_precise} simplifies to the quadratic control cost:
\begin{equation}\label{kl_control_cost}
\epsilon\KL(\mathrm{P}^u \| \mathrm{R}^{\epsilon,\mu_0}) = \mathbb{E}_{\mathrm{P}^u} \left[ \frac12 \int_0^{t_f} \|u(t,  X_t)\|^2 dt \right],
\end{equation}
showing that the Schr\"odinger bridge problem (SBP)~\eqref{sbp_precise} is equivalent to the stochastic control problem~\eqref{eq: objective}-\eqref{controlled_dynamics} on $(\R^d,\mathcal{D},g^{SR})$ with given initial and terminal distributions. Furthermore, from~\eqref{kl_control_cost}, we have that finite KL-divergence is equivalent to finite expected drift energy. Minimizing $\KL$ in~\eqref{sbp_precise} is equivalent to finding the minimal-energy control~\eqref{eq: objective}-\eqref{controlled_dynamics} to steer $\mu_0$ to $\mu_f$.  

The SBP  over a sub-Riemannian manifold $(\R^d,\mathcal{D},g^{SR})$ %\margin{we keep saying sub-R manifold, but we never say what we mean by this. We either should state what this means, or just say subject to the dynamics, which obviously is non-holonomic}
in~\eqref{sbp_precise} which serves as a stochastic counterpart to~\eqref{eq: expectation_total_control_effort}--\eqref{eq: Num_boundary condition} deviates from the literature. More precisely, in the classical setting (e.g. see~\cite{dai1991stochastic}) the reference process $dX_t = b(X_t)dt + g(X_t)dW_t$ has a uniformly elliptic (ue) diffusion matrix 
\[\xi^\top G(x) \xi \ge \lambda \|\xi\|^2, \quad\text{for any $\xi,x\in\R^d$} 
\]
where $G=gg^\top$ and for some constant $\lambda>0$. This is the situation in Figure~\ref{fig:diffusion_comparison_three}(c). Consequently its law $\mathrm{R}_{\mathrm{ue}}^{\epsilon}\in\mathcal P(\Omega)$ has full support on $\Omega$ and possesses a strictly positive smooth density. In contrast, $\mathrm{R}^{\epsilon,\mu_0}$ in~\eqref{eq: law of unreflected process} can be supported on the thin subset $\{X \in \Omega : dX_t \in \mathcal{D}_{X_t} \text{ for a.e. } t\}$, a proper, null-measure set for $\mathcal{R}^{\epsilon}_{ue}$, reflecting the loss of ellipticity. This is the situation in Figure~\ref{fig:diffusion_comparison_three}(b) and the SBP formulated in~\eqref{sbp_precise} will be ill-posed.  We show that the driftless characterization in~\eqref{uncontrolled_dynamic_constraint} is right setup to guarantee that the reference path space measure $\mathrm{R}^{\epsilon,\mu_0}$ has full support and admit a strictly positive transition density, which is necessary to characterize the solution of the SBP problem in~\eqref{sbp_precise}. This is the situation in Figure~\ref{fig:diffusion_comparison_three}(a).

To establish the existence of a unique solution of~\eqref{sbp_precise}, the following result will be useful.
%\margin{I cannot figure this out, but the lemma environment is giving us an error. Why?}
\begin{lemma} The infinitesimal generator 
\begin{equation}\label{eq:Ito-generator}
\mathcal L_\epsilon 
= b_\epsilon\cdot\nabla_x +\frac{\epsilon}{2}\mathrm{tr}\!\big(G\nabla_x^2\big)
\end{equation}
associated to~\eqref{uncontrolled_dynamic_constraint} can be expressed as the sum of square 
\begin{equation}\label{eq:sum-of-squares}
  \mathcal L_{\epsilon}= \tfrac12\sum_{i=1}^m V_i^2.
\end{equation}
where
\begin{equation}\label{eq: trans_vec_field}
V_i:=\sqrt{\epsilon}  g_i\cdot\nabla_x,\quad i=1,\dots,m.
\end{equation}
\end{lemma}
\begin{proof}
For a smooth test function $u=u(t,x)$, we have that
\[
V_i u
= \sqrt{\epsilon}\sum_{k=1}^d g^{k}_i\partial_{ x_k}u.
\]
Applying $V_i$ a second time and using the product rule,
\[
\begin{aligned}
V_i^2 u
&= V_i\left( \sqrt{\epsilon}\sum_{\ell=1}^d g^{\ell}_i\partial_{ x_\ell} u \right)
= \sqrt{\epsilon}\sum_{k=1}^d g^{k}_i\partial_{ x_k}
\left( \sqrt{\epsilon}\sum_{\ell=1}^d g^{\ell}_i\partial_{ x_\ell} u \right)\\
&= \epsilon\sum_{k,\ell=1}^d g^{k}_i(\partial_{ x_k} g^{\ell}_i)\partial_{ x_\ell} u+\epsilon\sum_{k,\ell=1}^d g^{k}_ig^{\ell}_i\partial_{ x_k  x_\ell}^2 u.
\end{aligned}
\]
Summing over $i=1,\dots,m$ and multiplying by $\tfrac12$ gives
\[
\frac12\sum_{i=1}^m V_i^2 u
= \frac{\epsilon}{2}\sum_{k,\ell=1}^d
\Big(\sum_{i=1}^m g^{k}_i g^{\ell}_i\Big)\partial_{ x_k  x_\ell}^2 u
+\frac{\epsilon}{2}\sum_{\ell=1}^d
\Big(\sum_{i=1}^m\sum_{k=1}^d g^{k}_i\partial_{ x_k} g^{\ell}_i\Big)\partial_{ x_\ell} u.
\]
Recognizing the diffusion matrix $G:=g g^\top$ with entries
$G_{k\ell}=\sum_{i=1}^m g^{k}_i g^{\ell}_i$,
we can rewrite the previous identity compactly as
\begin{equation}\label{eq:sum-of-squares-expanded}
\frac12\sum_{i=1}^m V_i^2 u
=\frac{\epsilon}{2}\mathrm{tr}\big(G\nabla_x^2 u\big)
+\frac{\epsilon}{2}\Big(\sum_{i=1}^m \nabla_{g_i}g_i\Big)\cdot\nabla_x u=\mathcal Lu.
\end{equation} 
\end{proof}

\begin{proposition}\label{lem:kernel}
The reference measure $\mathrm{R}^{\epsilon,\mu_0}\in\mathcal P(\Omega)$ in~\eqref{eq: law of unreflected process} associated to the uncontrolled SDE in~\eqref{uncontrolled_dynamic_constraint} 
admits a smooth positive transition density function. 
\end{proposition}
\begin{proof}
Using the Lie-algebra isomorphism $\Phi:v\mapsto v\cdot\nabla$, we have that
\[
\mathrm{Lie}(g_1,\dots,g_m)(x)= \mathrm{Lie}(V_1,\dots,V_m)(x)
\qquad \forall x\in\mathbb{R}^d,
\]
where $V_i$ are defined in~\eqref{eq: trans_vec_field}. Therefore, under Assumption~\eqref{ass:Controllability condition} and following from~\cite[Theorem~1.1]{Hormander1967} , we have that the operator $\partial_t-\mathcal L_{\epsilon}$, where $\mathcal L_{\epsilon}$ is in~\eqref{eq:sum-of-squares} is hypoelliptic. This implies that the PDE
\[
\partial_tu=\mathcal L_{\epsilon}^*u,\qquad  \forall s<t\qquad\text{and}\quad u(s)=\delta_{x},
\]
where $\delta_{x}$ is the Dirac distribution, admits a smooth transition density $p_{t-s,\epsilon}(x,y)$, for all $(x,y)\in\mathbb{R}^d\times\mathbb{R}^d$ and $t\geq s\geq 0$.   In our driftless case in~\eqref{eq:sum-of-squares}, controllable points can be reached using submersive controls~\cite[pp.~397]{arous1991decroissance} (controls that ensure the end-point map is a submersion). Hence under Assumption~\ref{ass:Controllability condition}, it follows that $p_{t,\epsilon}(x,y) > 0$ for all $x,y\in\R^d$ (see~\cite[Main Theorem]{leandre2005positivity}). 
%or~\cite{arous1991decroissance}.
\end{proof} 
Now that feasibility of~\eqref{sbp_precise} has been established, we proceed  to establish well-posedness of the SBP in~\eqref{eq: objective}-\eqref{controlled_dynamics} under degeneracy of the diffusion coefficient. 


\begin{theorem}\label{thm:properness}
The control problem~\eqref{eq: objective}-\eqref{controlled_dynamics} is proper: there exists an admissible $u$ that steers~\eqref{controlled_dynamics} from $ X_0^u\sim\mu_0$ to $ X_{t_f}^u\sim\mu_{t_f}$ and has finite energy in~\eqref{eq: objective}.% whenever $\rho_0,\rho_f$ $\int_{\R^d}\rho_f(y)\log(\rho_f(y))dy<\infty$.
\end{theorem}
\begin{proof}
We construct an admissible flow for~\eqref{controlled_dynamics} using Schr\"odinger bridge techniques. Let $\{ X_t\}_{0\leq t\leq t_f}$ satisfy the SDE in~\eqref{uncontrolled_dynamic_constraint} and $\mathcal F_{t}=\sigma( X_t)$ be the \emph{canonical filtration}, for all $t\in[0,t_f]$. 

% By strict positivity of the transition density function $p$ in~\eqref{eq:unref-rho0toTf},
For admissibility, from Proposition~\ref{lem:kernel}, following~\cite{Fortet1940}, there exist strictly positive functions $\map{\phi_0,\phi_f}{\mathbb R^d}{(0,\infty)}$ that satisfies  the Schr\"odinger system
\begin{equation}\label{eq:Sch-multplication_system}
\phi_0(x)(\mathscr P_{0,t_f} \phi_f)(x)=\rho_0(x)\qquad\text{and}\qquad (\mathscr P_{0,t_f}^* \phi_0)(y)\phi_f(y)=\rho_f(y),
\end{equation}
where 
\[
(\mathscr P_{s,t}\phi_f)(x):=\int \phi_f(y)p_{t-s,\epsilon}(x,y)dy\quad\text{and}\quad (\mathscr P_{s,t}^* \phi_0)(y):=\int \phi_0(x)p_{t-s,\epsilon}(x,y)dx.
\]
Define the measure on $\mathbb{R}^d\times \mathbb{R}^d$ as
\begin{equation}\label{eq: abs_distribution}
\frac{d\mathrm{Q}}{d\mathrm{R}^{\epsilon,\mu_0}}\Big|_{\mathcal F_{t_f}}:=\phi_0( x_0)\phi_f( x_{t_f}).    
\end{equation}
Then, using~\eqref{eq:Sch-multplication_system} and since
\[
\int_{(\mathbb{R}^d)^2}d\mathrm{Q}=\mathbb E_{\mathrm{R}^{\epsilon,\mu_0}}[\phi_0( x_0)\phi_f( x_{t_f})]
=\int_{\mathbb{R}^d} \phi_0(x)(\mathscr P_{0,t_f} \phi_f)(x)dx
=\int_{\mathbb{R}^d} \rho_0(x)dx=1,
\]
we have that $\mathrm{Q}$  is a probability measure. Furthermore, using~\eqref{eq:Sch-multplication_system}, we have that
\begin{equation*}
 \mathbb E_{\mathrm{Q}}[\psi( x_0)]
=\mathbb E_{\mathrm{R}^{\epsilon,\mu_0}} \big[\psi( x_0)\phi_0( x_0)\phi_f( x_{t_f})\big]
=\int \psi(x)\phi_0(x)(\mathscr P_{0,t_f} \phi_f)(x)dx =\int \psi(x)\rho_0(x)dx   
\end{equation*}
and
\begin{equation*}
\mathbb E_{\mathrm{Q}}[\psi( x_{t_f})]
=\mathbb E_{\mathrm{R}^{\epsilon,\mu_0}} \big[\psi( x_{t_f})\phi_0( x_0)\phi_f( x_{t_f})\big]
 =\int \psi(y)(\mathscr P_{0,t_f}^* \phi_0)(y)\phi_f(y)dy
 =\int \psi(y)\rho_f(y)dy    
\end{equation*}
hold for all bounded $\psi$. Therefore, 
\begin{equation}\label{eq: marginal_constraint}
( \mathrm{ev}_0)_{\#}\mathrm{Q}=\mu_0\qquad\text{and}\qquad ( \mathrm{ev}_{t_f})_{\#}\mathrm{Q}=\mu_{t_f}.    
\end{equation}
Restricting~\eqref{eq: abs_distribution} to a smaller $\sigma$-algebra $\mathcal F_{t}=\sigma( X_t)$ we have that 
\begin{equation}\label{eq:cond-step}
\frac{d\mathrm{Q}}{d\mathrm{R}^{\epsilon,\mu_0}}\Big|_{\mathcal F_{t}}
= \mathbb E_{\mathrm{R}^{\epsilon,\mu_0}} \big[\phi_0( x_0)\phi_f( x_{t_f})\ \big|\  X_t\big].   
\end{equation}
%(see, e.g., \cite[Th.\ 5.2]{Kallenberg2002} or the exercises around conditional expectation in \cite{Durrett2019,Billingsley1995,Williams1991}).
Using the Markov property (under $\mathrm{R}^{\epsilon,\mu_0}$, the past and future are conditionally
independent given $X_t$), the conditional expectation in \eqref{eq:cond-step} factorizes to
\[
\mathbb E_{\mathrm{R}^{\epsilon,\mu_0}} \big[\phi_0( x_0)\phi_f( x_{t_f})\ \big|\  X_t\big]
= \widehat\varphi(t, X_t)\varphi(t, X_t)
\]
where
\begin{equation}\label{eq: conditional expectation}
%\begin{aligned}
\widehat\varphi(t,x):=\mathbb E_{\mathrm{R}^{\epsilon,\mu_0}} \big[\phi_0( x_0)\mid  X_t=x\big] \quad\text{and}\quad%=(\mathscr P_{0,t}^* \phi_0)(x),\\[3pt]
\varphi(t,x):=\mathbb E_{\mathrm{R}^{\epsilon,\mu_0}} \big[\phi_f( x_{t_f})\mid  X_t=x\big].%=(\mathscr P_{t,t_f}\phi_f)(x).
%\end{aligned}    
\end{equation}
Therefore, the density process of~\eqref{eq: abs_distribution} is 
\[
Z_t=\frac{d\mathrm{Q}}{d\mathrm{R}^{\epsilon,\mu_0}}\Big|_{\mathcal F_t}=\widehat\varphi(t, X_t)\varphi(t, X_t).
\]
By the Doob/L\'eonard generalized $h$-transform~ \cite{Doob1984,Leonard2011StochasticH,Leonard2014Survey}, we have that $\mathrm{Q}$ is Markov with generator
\[
\mathcal L_\epsilon^{\mathrm{Q}} f=b_\epsilon\cdot\nabla f+\epsilon G\nabla\log\varphi\cdot\nabla f+\tfrac{\epsilon}{2}\mathrm{tr}(G\nabla^2 f).
\]
Equivalently, we have that $\mathrm{Q}$ is the distribution of the process
\begin{equation}\label{eq:unref-tiltedSDE}
d x_t=(b_{\epsilon}(X_t)+\epsilon  G(X_t)  \nabla\log\varphi(t, X_t)) dt+\sqrt{\epsilon}  g(X_t)  dW_t^{\mathrm{Q}}
\end{equation}
that satisfies the marginal constraint in~\eqref{eq: marginal_constraint}.
One can check that the above is the controlled SDE in~\eqref{controlled_dynamics} with feedback
\begin{equation}\label{eq: admissible control}
u(t,x)=\epsilon  g(x)^\top\nabla\log\varphi(t,x),
\end{equation}
which concludes admissibility.

Consider the reference endpoint coupling,
\[
r(dx,dy)=(\mathrm{ev}_0,\mathrm{ev}_{t_f})_{\#}\mathrm{R}^{\epsilon,\mu_0}=p_{t_f,\epsilon}(x,y)\rho_0(x)dxdy.%K(x,dy)\mu_0(dx),
\]
%where $K(x,dy)=p_{t_f,\epsilon}(x,y)dy$ is the final time transition. 
By disintegration of the relative entropy, we have that 
\begin{equation}\label{eq:KL-disintegration}
\KL(\mathrm{Q} \| \mathrm{R}^{\epsilon,\mu_0})
\;=\;
\KL(\pi \| r)
\;+\;
\int_{\R^d\times\R^d}
\KL\!\big(\mathrm{Q}^{x,y} \| \R^{\epsilon,x,y}\big) \pi(dx,dy),
\end{equation}
where $\pi\in \Pi(\mu_0,\mu_f)$ in~\eqref{eq: transport_plans} and $\mathrm{R}^{\epsilon,x,y}( \cdot )=\mathrm{R}^{\epsilon,\mu_0}( \cdot \mid \mathrm{ev}_0(\cdot)=x,\mathrm{ev}_{t_f}(\cdot)=y)$ is the bridge of~\eqref{uncontrolled_dynamic_constraint} and similarly for $\mathrm{Q}^{x,y}$ in~\eqref{eq:unref-tiltedSDE}. From~\eqref{eq: abs_distribution}, since 
\begin{align*}
\E_\mathrm{Q}[F\mid \mathrm{ev}_0(\omega)=x,\mathrm{ev}_{t_f}(\omega)=y]%=&\frac{\E_{\mathrm{R}^{\epsilon,\mu_0}}[F\phi_0\phi_f\mid \mathrm{ev}_0(\omega)=x,\mathrm{ev}_{t_f}(\omega)=y]}{\E_{\mathrm{R}^{\epsilon,\mu_0}}[\phi_0\phi_f\mid \mathrm{ev}_0(\omega)=x,\mathrm{ev}_{t_f}(\omega)=y]}\\
=& ~\E_{\mathrm{R}^{\epsilon,\mu_0}}[F\mid \mathrm{ev}_0(\omega)=x,\mathrm{ev}_{t_f}(\omega)=y]
\end{align*}
%\margin{I guess the last step is clear to reader? After conditioning on end points, $\phi_0$ and $ \phi_f $ are constants and so they cancel. Not sure this needs to be mentioned. Okay if not}
holds, for every bounded measurable \(F:\Omega\to\R\),  we have that 
\[
\mathrm{Q}^{x,y} = \mathrm{R}^{\epsilon,x,y}, \quad\text{for}\qquad r\text{-a.e. }(x,y).
\]
Therefore,
\[
\min_{\mathrm{P}^u \in \mathcal C_{\mathrm{SB}}} \KL(\mathrm{P}^u \| \mathrm{R}^{\epsilon,\mu_0})=\min_{\pi\in\Pi(\mu_0,\mu_f)} \KL(\pi \| r).
\]
Since $\mu_0\otimes\mu_f\in \Pi(\mu_0,\mu_f)$ this gives an upper bound
\begin{multline}\label{eq: finite_energy_bound}
\min_{\mathrm{P}^u \in \mathcal C_{\mathrm{SB}}} \KL(\mathrm{P}^u \| \mathrm{R}^{\epsilon,\mu_0})\leq \KL(\mu_0\otimes\mu_f \| r)\\=\int_{\R^d}\rho_f(y)\log(\rho_f(y))dy-\int_{\R^d\times\R^d}\log(p_{t_f,\epsilon}(x,y))\rho_0(x)\rho_f(y)dxdy
\end{multline}
where the last inequality is finite due to the fact that $\rho_0,\rho_f$ have compact support (or having exponential tail) on $\R^d$ and $p_{t_f,\epsilon}(\cdot,\cdot)$ is positive and smooth. This completes the proof.
\end{proof}
Although the skeleton of the proof is similar to~\cite{Leonard2014Survey}, the degeneracy of $G = gg^\top$ in~\eqref{eq:unref-tiltedSDE} forces the control~\eqref{eq: admissible control} to lie in $\operatorname{span}\{g_1,\dots,g_m\}$. Moreover, the finite-energy bound in~\eqref{eq: finite_energy_bound} relies on properties of the hypoelliptic heat-kernel, which in turn depend on the sub-Riemannian distance on~\eqref{eq: distribution}.


\section{Convex Eulerian Formulation}\label{sec:Eulerian}
To facilitate the proof of convergence from the horizontal Schr\"odinger bridge problem in~\eqref{sbp_precise} to the sub-Riemannian optimal-transport problem~\eqref{eq: expectation_total_control_effort}-\eqref{eq: Num_boundary condition}, we work in the Eulerian formulation.

For a fixed $\epsilon>0$, since the reference measure $\mathrm{R}^{\epsilon,\mu_0}\in\mathcal P(\Omega)$ of~\eqref{uncontrolled_dynamic_constraint} admits a density function and the probability distribution $\mathrm{P}^u$  of the process $\{ X_t^u\}_{0\leq t\leq t_f}\in \mathbb{R}^d$ in~\eqref{controlled_dynamics} is absolutely continuous with respect to $\mathrm{R}^{\epsilon,\mu_0}$, we have that $\mathrm{P}^{u_{\epsilon}}$  admits a density function which we denote as $\rho_{\epsilon}$. Then,
\begin{equation}\label{eq: time_varying density}
 \rho_{\epsilon}(t,\cdot) \in L^1\big(\mathbb{R}^d\big), \quad \rho_{\epsilon} \geq 0,\quad \int_{\R^d} \rho_{\epsilon}(t,x)  dx = 1 \quad \text{for a.e. } t
\end{equation}
and the cost functional in~\eqref{eq: objective} becomes %the expectation (with respect to the law $\rho(t,x)$ of the state) of the control cost for the original optimal control problem~\eqref{controlled_dynamics}:
\begin{equation}\label{eq: function_objective}
\mathbb{E} \Big[ \frac12 \int_0^{t_f} \tfrac{1}{2} \| u_{\epsilon}(t,  X_t^u) \|^2 \Big] =  \frac12 \int_0^{t_f}\int_{\mathbb{R}^d} \frac12 \| u_{\epsilon}(t,x) \|^2    \rho_{\epsilon}(t,x)    dx.
\end{equation}
Motivated by~\cite{BenamouBrenier2000}, let $m$ be absolutely continuous w.r.t. $\rho$, meaning
\begin{equation}\label{eq: momentum_pointwise}
m_{\epsilon}(t,x) = u_{\epsilon}(t,x)\rho_{\epsilon}(t,x)  \quad \text{for some } u_{\epsilon} \in L^2(\rho_{\epsilon}; \mathbb{R}^m).
\end{equation}
Then, from~\eqref{eq: function_objective}, we have
\[
\frac12\int_0^{t_f}\!\int_{\R^d}\|u_{\epsilon}(t,x)\|^2 \rho_{\epsilon}(t,x) dxdt
=
\int_Q \frac{\|m_{\epsilon}(t,x)\|^2}{2\rho_{\epsilon}(t,x)} dt dx,
\]
where $Q:=(0,t_f)\times\R^d$ and the function pair $(\rho_{\epsilon},m_{\epsilon})$ solves the controlled Fokker--Planck equation
\begin{subequations}
\begin{align}
   \partial_t \rho_{\epsilon}
+ \nabla_x \cdot \left(b_{\epsilon}\rho_{\epsilon}+gm_{\epsilon}\right)
- \frac{\epsilon}{2} \sum_{i,j=1}^d 
\partial_{x_i x_j}^2 \Big(G_{ij} \rho_{\epsilon} \Big)
&= 0\quad\text{in }\quad \bigl(C_c^\infty(Q)\bigr)^{*}\label{eq: fb_constraint}\\
    \rho_{\epsilon}(0, \cdot) = \rho_0, 
    \qquad 
    \rho_{\epsilon}(t_f, \cdot) &= \rho_{t_f},\quad\quad\text{in\quad $\mathcal P(\R^d)$.} 
    \label{eq:endpoint_constraints}
\end{align}
where $\bigl(C_c^\infty(Q)\bigr)^{*}$ is the dual of  the space of compactly supported smooth functions $\bigl(C_c^\infty(Q)\bigr)$.  
\end{subequations}
The Eulerian problem corresponding to~\eqref{eq: objective}-\eqref{controlled_dynamics} reads:
\begin{equation}\label{eq: convexified_stochastic}
\inf_{\substack{(\rho_{\epsilon}, m_{\epsilon})\\ \text{\eqref{eq: fb_constraint}-\eqref{eq:endpoint_constraints}}}} \mathcal{J}(\rho_{\epsilon},m_{\epsilon}):=
\int_0^{t_f} \int_{\mathbb R^d}
\frac{\| m_{\epsilon}(t,x) \|_2^2}{2   \rho_{\epsilon}(t,x)}   dx   dt.
%\quad \text{subject to~\eqref{eq: fb_constraint}-\eqref{eq:endpoint_constraints}.}
\end{equation}

One may regard $\boldsymbol\mu_{\epsilon}\in \mathcal P(Q)$ as a nonnegative real-valued probability measure on $Q$ with density $\rho_{\epsilon}$ in~\eqref{eq: time_varying density}, that is
\[
\boldsymbol\mu_{\epsilon}(dt,dx):=\rho_{\epsilon}(t,x) dt dx,
\]
and $\mathbf m_{\epsilon} \in\mathcal{M}(Q;\R^m)$ as an $\R^m$-valued Radon measure with density $m_{\epsilon}$. The statement in~\eqref{eq: momentum_pointwise} precisely means: $\mathbf m_{\epsilon}\ll \boldsymbol\mu_{\epsilon}$ and $u_{\epsilon}=\frac{d\mathbf m_{\epsilon}}{d\boldsymbol\mu_{\epsilon}}\in L^2(\boldsymbol\mu_{\epsilon};\R^m)$. In that case, consider
\[
 \mathcal C: =\bigl\{ (\boldsymbol\mu_{\epsilon}, \mathbf m_{\epsilon})\in \mathcal{P}(\overline Q)\times \mathcal{M}(\overline Q) : \mathbf m_{\epsilon}\ll \boldsymbol\mu_{\epsilon},\ 
\frac{d\mathbf m_{\epsilon}}{d\boldsymbol\mu_{\epsilon}}\in L^2(\boldsymbol\mu_{\epsilon};\R^m) \bigr\}.
\]
Then we can express \eqref{eq: convexified_stochastic} as the optimization problem,
\begin{equation}\label{measure_theoretic_objective}
\inf_{(\boldsymbol\mu_{\epsilon}, \mathbf m_{\epsilon})\in\mathcal C} 
\mathcal{J}(\boldsymbol\mu_{\epsilon}, \mathbf m_{\epsilon}):=
\int_Q \frac12\Big\|\frac{d\mathbf m_{\epsilon}(t,x)}{d\boldsymbol\mu_{\epsilon}(t,x)}\Big\|^2 d\boldsymbol\mu_{\epsilon}.%(dt,dx)
\end{equation}
%\margin{Very important one: I have been confused throughout by when we use $ \boldsymbol\mu_{\epsilon}$ vs $\boldsymbol\mu_\epsilon$ to refer to the solution of the above. You have $ \mathcal{J}$ in here, and so perhaps it made sense to use $\boldsymbol\mu_\epsilon$. Now, if you look at say (4.3), there you use $\boldsymbol\mu_\epsilon$. This is one of those changes that if done, needs to be done very carefully, and consistently. }
subject to
\begin{subequations}
\begin{align}
   \partial_t \boldsymbol\mu_{\epsilon}
+ \nabla_x \cdot \left(b_{\epsilon}\boldsymbol\mu_{\epsilon}+g\mathbf m_{\epsilon}\right)
- \frac{\epsilon}{2} \sum_{i,j=1}^d 
\partial_{ x_i  x_j}^2 (G_{ij} \boldsymbol\mu_{\epsilon} )
&= 0\quad\text{in } \bigl(C_c^\infty(Q)\bigr)^{*},\label{eq: fb_measure_constraint}\\
    \boldsymbol\mu_{\epsilon}|_{t=0} = \mu_0, 
    \qquad 
    \boldsymbol\mu_{\epsilon}|_{t=t_f} &= \mu_f,\quad\text{in } \mathcal P(\R^d).\label{eq:endpoint_measure_constraints}
\end{align}
\end{subequations}
Since the constrained 
%Fokker-Planck equation in~\eqref{measure_theoretic_objective}-\eqref{eq:endpoint_measure_constraints} 
is linear in $(\boldsymbol{\mu}_\epsilon,\mathbf{m}_\epsilon)$ and $(\boldsymbol{\mu}_\epsilon,\mathbf{m}_\epsilon)\mapsto \mathcal{J}(\boldsymbol{\mu}_\epsilon,\mathbf{m}_\epsilon)$ is jointly strictly convex (see~\cite{BenamouBrenier2000}), the Schr\"odinger bridge analysis in Theorem~\eqref{thm:properness} shows that the problem~\eqref{eq: convexified_stochastic} is proper. 
We state the existence and uniqueness result without proof.
\begin{theorem}\label{thm: existence and uniqueness}
Problem~\eqref{measure_theoretic_objective} subject to~\eqref{eq: fb_measure_constraint}-\eqref{eq:endpoint_measure_constraints} admits a unique minimizer $(\boldsymbol\mu_{\epsilon}^*, \mathbf m_{\epsilon}^*)$.  
\end{theorem} 

We proceed to show that the unique measure minimizer $(\boldsymbol{\mu}_\epsilon^*, \mathbf{m}_\epsilon^*)$ are actually given by $d\boldsymbol{\mu}_\epsilon^* = \rho_\epsilon^* dt dx$, $d\mathbf{m}_\epsilon^* = m_{\epsilon}^\star dt dx$, where $(\rho_\epsilon^*,m_{\epsilon}^\star)$ are minimizers of Problem~\eqref{eq: convexified_stochastic}.
\begin{proposition}\label{prop: Strong duality and attainment}
Let 
\begin{equation}\label{eq: dual_set_thm}
 \Lambda := \Bigg\{ \lambda \in C^{1,2}([0,t_f]\times\mathbb{R}^d) \;\Bigg|\; \partial_t\lambda + b_\epsilon\cdot\nabla\lambda + \frac12\langle \nabla\lambda, G\nabla\lambda\rangle + \frac{\epsilon}{2}\operatorname{tr}(G\nabla^2\lambda) \le 0 \quad \text{on } (0,t_f)\times\mathbb{R}^d \Bigg\}   
\end{equation}
and define the dual functional $\map{\mathcal D}{\Lambda}{\{\R\cup-\infty\}}$ as
\begin{equation}\label{eq:dual-func-thm}
\mathcal D(\lambda):=\int_{\R^d}\lambda(t_f,x)\rho_f(x) dx
-\int_{\R^d}\lambda(0,x)\rho_0(x) dx.
\end{equation}

Assume that there exist strictly positive functions $(\varphi_\epsilon,\hat\varphi_\epsilon)$ that solve: %the Schr\"odinger system
\begin{subequations}\label{eq:Schr-system-thm}
\begin{align}
\partial_t\varphi_\epsilon + \mathcal L_\epsilon \varphi_\epsilon &=0,\label{eq:Schr-fwd}\\
\partial_t\hat\varphi_\epsilon - \mathcal L_\epsilon^* \hat\varphi_\epsilon &=0,\label{eq:Schr-bwd}\\
\varphi_\epsilon(0,\cdot)\hat\varphi_\epsilon(0,\cdot)=\rho_0,\qquad
\varphi_\epsilon(t_f,\cdot)\hat\varphi_\epsilon(t_f,\cdot)&=\rho_f,\label{eq:Schr-bc}
\end{align}
\end{subequations}
where $\mathcal L_\epsilon$ is in~\eqref{eq:Ito-generator}. Then strong duality holds: 
\[ 
\inf_{\substack{(\rho_{\epsilon}, m_{\epsilon}) \\ \text{\eqref{eq: fb_constraint}-\eqref{eq:endpoint_constraints}}}} 
\mathcal{J}(\rho_{\epsilon}, m_{\epsilon}) = \sup_{\lambda_{\epsilon}\in\Lambda}\mathcal D(\lambda_{\epsilon})
\]
and the supremum is attained at
\begin{equation}\label{eq:optimal-dual}
\lambda_\epsilon^\star:=\epsilon\log\varphi_\epsilon.
\end{equation}
Moreover, the primal infimum is uniquely attained at
\begin{align}\label{eq:optimal-primal}
&\rho_\epsilon^\star:=\varphi_\epsilon\hat\varphi_\epsilon\in C^\infty((0,t_f)\times\R^d),\nonumber\\\
&m_\epsilon^\star:=\rho_\epsilon^\star g^\top\nabla\lambda_\epsilon^\star = \epsilon \hat\varphi_\epsilon g^\top\nabla\varphi_\epsilon\in C^\infty((0,t_f)\times\R^d).
\end{align}    
\end{proposition}
\begin{proof}
We establish the result in three steps: weak duality, feasibility of the constructed pair, and equality verifying attainment.

For weak duality, for any primal feasible $(\rho,m)$ and dual feasible $\lambda\in\Lambda$, we test the Fokker-Planck constraint in~\eqref{eq: fb_constraint}-\eqref{eq:endpoint_constraints} against $\lambda$ (approximating by compactly supported test functions if necessary). Hence integration by parts yields
\begin{align*}
\mathcal D(\lambda_{\epsilon})
&=\int_0^{t_f}\!\!\int_{\R^d}
\Bigl[\rho_{\epsilon}\bigl(\partial_t\lambda_{\epsilon}+b_\epsilon\cdot\nabla\lambda_{\epsilon}
+\tfrac{\epsilon}{2}\mathrm{tr}(G\nabla^2\lambda_{\epsilon})\bigr)
+ m\cdot(g^\top\nabla\lambda_{\epsilon})\Bigr] dx dt.
\end{align*}
Since $\lambda_{\epsilon}\in\Lambda$ in~\eqref{eq: dual_set_thm}, we get that
\[
\partial_t\lambda_{\epsilon}+b_\epsilon\cdot\nabla\lambda_{\epsilon}
+\tfrac{\epsilon}{2}\mathrm{tr}(G\nabla^2\lambda_{\epsilon})
\le -\tfrac12\langle\nabla\lambda_{\epsilon},G\nabla\lambda_{\epsilon}\rangle
=-\tfrac12\|g^\top\nabla\lambda_{\epsilon}\|^2,
\]
and hence
\[
\mathcal D(\lambda_{\epsilon})\le \int_0^{t_f}\!\!\int_{\R^d}
\Bigl(m_{\epsilon}\cdot(g^\top\nabla\lambda_{\epsilon})-\tfrac{\rho_{\epsilon}}{2}\|g^\top\nabla\lambda_{\epsilon}\|^2\Bigr) dx dt.
\]
Applying Young's inequality pointwise gives
\begin{equation}\label{eq: Young's inequality}
m_{\epsilon}\cdot(g^\top\nabla\lambda_{\epsilon})=\left(\frac{m_{\epsilon}}{\sqrt{\rho_{\epsilon}}}\right)\left(\sqrt{\rho_{\epsilon}}(g^\top\nabla\lambda_{\epsilon}\right)
\le \frac{\|m_{\epsilon}\|^2}{2\rho_{\epsilon}}+\frac{\rho_{\epsilon}}{2}\|g^\top\nabla\lambda_{\epsilon}\|^2,
\end{equation}
and hence we obtain 
\[
\mathcal D(\lambda_{\epsilon})\le \mathcal J(\rho_{\epsilon},m_{\epsilon})
\]
for all $\lambda_{\epsilon}\in\Lambda$ in~\eqref{eq: dual_set_thm} and the pairs $(\rho_{\epsilon},m_{\epsilon})$ in~\eqref{eq: fb_constraint}-\eqref{eq:endpoint_constraints}. Thus 
\begin{equation}\label{eq: weak duality} 
\inf_{\substack{(\boldsymbol\mu_{\epsilon}, \mathbf m_{\epsilon})\in \mathcal C \\ \text{\eqref{eq: fb_constraint}-\eqref{eq:endpoint_constraints}}}} 
\mathcal{J}(\rho_{\epsilon},m_{\epsilon}) \geq \sup_{\lambda_{\epsilon}\in\Lambda}\mathcal D(\lambda_{\epsilon}).
\end{equation}
%\margin{should $ \rho_{\epsilon}^\star$ be $\rho_\epsilon^\star$}
For feasibility and equality in the HJB, we consider the function tuple $(\lambda_{\epsilon}^\star,\rho_{\epsilon}^\star,m_{\epsilon}^\star)$ in \eqref{eq:optimal-dual}--\eqref{eq:optimal-primal}. From \eqref{eq:Schr-bc}, $\rho_{\epsilon}^\star$ satisfies the endpoint conditions~\eqref{eq:endpoint_constraints}. A direct computation using $\lambda_{\epsilon}^\star=\epsilon\log\varphi_{\epsilon}$ and \eqref{eq:Schr-fwd} gives
\begin{multline*}
\partial_t\lambda^\star+b_\epsilon\cdot\nabla\lambda^\star
+\tfrac12\langle\nabla\lambda^\star,G\nabla\lambda^\star\rangle
+\tfrac{\epsilon}{2}\tr(G\nabla^2\lambda^\star) \\
= \epsilon\frac{\partial_t\varphi_{\epsilon}+b_\epsilon\cdot\nabla\varphi_{\epsilon}+\tfrac{\epsilon}{2}\tr(G\nabla^2\varphi_{\epsilon})}{\varphi_{\epsilon}}
-\tfrac{\epsilon^2}{2}\frac{\langle\nabla\varphi_{\epsilon},G\nabla\varphi_{\epsilon}\rangle}{\varphi_{\epsilon}^2}
+\tfrac{\epsilon^2}{2}\frac{\langle\nabla\varphi_{\epsilon},G\nabla\varphi_{\epsilon}\rangle}{\varphi_{\epsilon}^2}
\\= \epsilon\frac{\partial_t\varphi_{\epsilon}+\mathcal L_\epsilon\varphi_{\epsilon}}{\varphi_{\epsilon}}=0.
\end{multline*}
Thus $\lambda^\star\in\Lambda$, achieving the equality in \eqref{eq: dual_set_thm}. To verify primal feasibility, note that $m_{\epsilon}^\star=\rho_{\epsilon}^\star g^\top\nabla\lambda^\star$ by construction. Using \eqref{eq:Schr-system-thm}, the pair $(\rho_{\epsilon}^\star,m_{\epsilon}^\star)$ satisfies
\begin{align*}
\partial_t\rho_{\epsilon}^\star
&= \hat\varphi_{\epsilon}\,\partial_t\varphi_{\epsilon} + \varphi_{\epsilon}\,\partial_t\hat\varphi_{\epsilon}
= -\hat\varphi_{\epsilon}\,\mathcal L_\epsilon\varphi_{\epsilon} + \varphi_{\epsilon}\,\mathcal L_\epsilon^\ast\hat\varphi_{\epsilon} \\
&= -\hat\varphi_{\epsilon}\Big(b_\epsilon\cdot\nabla\varphi_{\epsilon} + \frac{\epsilon}{2}\tr(G\nabla^2\varphi_{\epsilon})\Big)
\;+\;\varphi_{\epsilon}\Big(-\nabla\!\cdot(b_\epsilon\hat\varphi_{\epsilon})+\frac{\epsilon}{2}\sum_{i,j=1}^d\partial_{ij}(G_{ij}\hat\varphi_{\epsilon})\Big) \\
&= -\nabla\!\cdot(b_\epsilon\,\rho_{\epsilon}^\star)
\;+\;\frac{\epsilon}{2}\sum_{i,j=1}^d \partial_{ij}\!\big(G_{ij}\rho_{\epsilon}^\star\big)
\;-\;\epsilon\,\nabla\!\cdot\!\big(\hat\varphi_{\epsilon}\,G\nabla\varphi_{\epsilon}\big).
\end{align*}
where the cross terms combine to yield precisely $\nabla\cdot(gm_{\epsilon}^\star) = \nabla\cdot(\epsilon\hat\varphi_{\epsilon} G\nabla\varphi_{\epsilon})$. Hence the pair $(\rho_{\epsilon}^\star,m_{\epsilon}^\star)$ solves~\eqref{eq: fb_constraint}. Furthermore, since
%\[
%\frac{\|m^*\|^2}{\rho_{\epsilon}^\star} = \frac{\epsilon^2 \hat{\varphi}^2 \|g^\top \nabla\varphi\|^2}{\varphi\hat{\varphi}} = \epsilon^2 \hat{\varphi} \frac{\|g^\top \nabla\varphi\|^2}{\varphi}
%\]
%and hence
%Using the identity $\partial_t\varphi + \mathcal{L}_\epsilon\varphi = 0$, one shows that
\[
\int_Q \frac{\|m_{\epsilon}^\star\|^2}{2\rho_{\epsilon}^\star} dxdt = \epsilon^2 \int_Q \hat{\varphi_{\epsilon}} \frac{\|g^\top \nabla\varphi_{\epsilon}\|^2}{2\varphi_{\epsilon}} dxdt= \epsilon^2\int_Q \hat{\varphi_{\epsilon}}  \varphi_{\epsilon} \frac{\|g^\top \nabla(\log\varphi_{\epsilon})\|^2}{2}   dxdt,
\]
%\margin{Where is the $1/2 $ coming from? I think there is a minor error here}
and
\[
\epsilon^2\int_Q \hat{\varphi_{\epsilon}}   \varphi_{\epsilon}   \frac{\|g^\top \nabla(\log\varphi_{\epsilon})\|^2}{2}   dxdt=\int_Q  \frac12\|u^*\|^2\rho_{\epsilon}^\star   dxdt
\]
from~\eqref{kl_control_cost} and the energy estimates in~\eqref{eq: finite_energy_bound} we conclude that
\[
\int_Q \frac{\|m_{\epsilon}^\star\|^2}{\rho_{\epsilon}^\star} dxdt<\infty.
\]

For attainment and strong duality, note that if $m_{\epsilon}^\star=\rho_{\epsilon}^\star g^\top\nabla\lambda_{\epsilon}^\star$, then the equality holds in Young's inequality in~\eqref{eq: Young's inequality} and equality holds in~\eqref{eq: weak duality}. Therefore
\[
\mathcal D(\lambda_{\epsilon}^\star) = \mathcal J(\rho_{\epsilon}^\star,m_{\epsilon}^\star).
\]
Combined with weak duality, this yields
\[
\mathcal D(\lambda^\star) = \mathcal J(\rho_{\epsilon}^\star,m_{\epsilon}^\star) = \inf\mathcal J = \sup_{\Lambda}\mathcal D,
\]
proving that $\lambda_{\epsilon}^\star$ attains the dual supremum, $(\rho_{\epsilon}^\star,m_{\epsilon}^\star)$ attains the primal infimum, and there is no duality gap.

For regularity, note that the solution for~\eqref{eq:Schr-fwd}-\eqref{eq:Schr-bc} is characterized as
\[
\varphi_\epsilon(t,\cdot) = e^{(t_f-t)\mathcal{L}_\epsilon}\varphi_{t_f}, \qquad 
\hat{\varphi}_\epsilon(t,\cdot) = e^{t\mathcal{L}_\epsilon^*}\hat{\varphi}_0,
\]
with $\hat{\varphi}_0 = \rho_0 / \varphi(0,\cdot)$ and $\varphi_{t_f}=\phi_f$ in~\eqref{eq:Sch-multplication_system},  where the forward and backward semigroups $P_t = e^{t\mathcal{L}_\epsilon^*}$ and $
Q_t = e^{t\mathcal{L}_\epsilon}$ are defined through the actions
\begin{equation}\label{eq: action_semigroup}
(P_t \varphi_0)(y) = \int_{\R^d} p_{t,\epsilon}(x,y) \varphi_0(x)   dx\quad\text{and}\quad (Q_t \varphi_f)(x) = \int_{\R^d} p_{t,\epsilon}(x,z) \varphi_f(z)   dz    
\end{equation}
respectively and $p_{t,\epsilon}(x,y)$ is the transition density function. Under Assumption~\ref{ass:Controllability condition}, since the kernels $p_{t,\epsilon}(x,y)$ are smooth and strictly positive for all $t>0$ we have that the functions $\varphi$ and $\hat{\varphi}$  inherit this smoothness and strict positivity. Hence $\rho_\epsilon^* = \varphi_\epsilon\hat{\varphi}_\epsilon \in C^\infty((0,t_f)\times\R^d)$ and $\rho_{\epsilon}^\star > 0$ everywhere. Moreover, $m_{\epsilon}^\star$ is smooth by the smoothness of $\varphi_\epsilon$, $\hat{\varphi}_\epsilon$, and $g$. The uniqueness follows from the strict convexity of $\mathcal{J}$ and the linearity of the constraint set.
\end{proof}
By the uniqueness in Theorem~\ref{thm: existence and uniqueness} and Proposition~\ref{prop: Strong duality and attainment}, we conclude here that the measure minimizer $(\boldsymbol{\mu}_{\epsilon}^*, \mathbf{m}_{\epsilon}^*)$ is actually given by $d\boldsymbol{\mu}_{\epsilon}^* = \rho_{\epsilon}^* dt dx$, $d\mathbf{m}_{\epsilon}^* = m_{\epsilon}^* dt dx$, where $(\rho_{\epsilon}^*,m_{\epsilon}^*)\in C^\infty((0,t_f)\times\R^d)\times C^\infty((0,t_f)\times\R^d;\R^d)$ are in~\eqref{eq:optimal-primal}. From this the optimal control is characterized as
\begin{equation}\label{eq:optimal-control}
u_\epsilon^\star:=\frac{m_{\epsilon}^\star}{\rho_{\epsilon}^\star}=g^\top\nabla\lambda_{\epsilon}^\star=\epsilon  g^\top\nabla\log\varphi_{\epsilon},
\end{equation}
where $\varphi_{\epsilon}$ is a strictly positive functions and together with $\hat\varphi_{\epsilon}$ solves the Schr\"odinger system~\eqref{eq:Schr-fwd}-\eqref{eq:Schr-bc}.


\section{From Schr\"odinger Bridge to Optimal Transport}\label{sec: Limiting Section}
In this section we analyze the zero-noise limit $\epsilon \downarrow 0$ of the entropic (horizontal) Schr\"odinger bridge and show that the sequence of solutions $(\boldsymbol{\mu}_\epsilon^*, \mathbf{m}_\epsilon^*)$ of Problem~\eqref{measure_theoretic_objective} converges, up to subsequences, to a solution of the sub-Riemannian Benamou--Brenier optimal transport problem for~\eqref{eq: expectation_total_control_effort}-\eqref{eq: Num_boundary condition}.
%In this section we analyze the zero-noise limit of the entropic (horizontal) Schr\"odinger bridge problem ~\eqref{eq: objective}-\eqref{controlled_dynamics}  and show that it converges (in the sense of $\Gamma$-convergence~\cite{dal2012introduction})to the sub-Riemannian Benamou--Brenier optimal transport problem~\eqref{eq: expectation_total_control_effort}-\eqref{eq: Num_boundary condition}. We proceed to the result in the Eulerian settings.
\begin{theorem}\label{thm:Gamma}
Consider the problem
\begin{equation}\label{measure_theoretic_objective2}
\inf_{(\boldsymbol\mu, \mathbf m)\in\mathcal C} 
\mathcal{J}(\boldsymbol\mu, \mathbf m):=
\int_Q \frac12\Big\|\frac{d\mathbf m(t,x)}{d\boldsymbol\mu(t,x)}\Big\|^2 d\boldsymbol\mu%(dt,dx)
\end{equation}
subject to the continuity equation
\begin{subequations}\label{eq:CE_measure_unref}
\begin{align}
\partial_t \boldsymbol\mu + \nabla_x \cdot \left(g\mathbf m\right) &= 0 \quad \text{in } \bigl(C_c^\infty(Q)\bigr)^*, \label{eq:CE_measure_constraint}\\
\boldsymbol\mu|_{t=0} = \mu_0, \qquad \boldsymbol\mu|_{t=t_f} &= \mu_f \quad \text{in } \mathcal{P}(\mathbb{R}^d). \label{eq:CE_measure_endpoint}
\end{align}
\end{subequations}
 Let $(\boldsymbol{\mu}_\epsilon^*, \mathbf{m}_\epsilon^*)$ be the unique minimizer of Problem~\eqref{measure_theoretic_objective} for each $\epsilon > 0$. Then along the convergent subsequence, we have that
        \begin{equation}\label{eq:aux1}
            \lim_{\epsilon \to 0} \mathcal{J}(\boldsymbol{\mu}_\epsilon^*, \mathbf{m}_\epsilon^*) = \mathcal{J}(\boldsymbol{\mu}^*, \mathbf{m}^*) = \min_{(\boldsymbol\mu,\mathbf m)\text{ satisfying \eqref{eq:CE_measure_constraint}--\eqref{eq:CE_measure_endpoint}}} \mathcal J(\boldsymbol\mu,\mathbf m).
        \end{equation}
\end{theorem}

\begin{proof}
 We show this in three steps; Compactness of the sequence of minimizers, the limit sequence satisfying continuity equation and optimality of the limit sequence through energy convergence. Since the unique minimizer $(\boldsymbol{\mu}_\epsilon^*, \mathbf{m}_\epsilon^*)$ of Problem~\eqref{measure_theoretic_objective} is characterized as $d\boldsymbol{\mu}_{\epsilon}^* = \rho_{\epsilon}^* dt dx$, $d\mathbf{m}_{\epsilon}^* = m_{\epsilon}^* dt dx$, where $(\rho_{\epsilon}^*,m_{\epsilon}^*)\in C^\infty((0,t_f)\times\R^d)\times C^\infty((0,t_f)\times\R^d;\R^d)$ are in~\eqref{eq:optimal-primal}, from~\eqref{kl_control_cost},~\eqref{eq: function_objective} we have that
\[
\int_Q \frac12\Big\|\frac{d\mathbf m_{\epsilon}^\star(t,x)}{d\boldsymbol\mu^*_\epsilon(t,x)}\Big\|^2 d\boldsymbol\mu^*_\epsilon=\int_Q  \frac{\|m_{\epsilon}^\star(t,x)\|^2}{2\rho_\epsilon^*(t,x)} dx\, dt =\epsilon\KL(\mathrm{P}^{u_\epsilon^\star} \| \mathrm{R}^{\epsilon,\mu_0})
\]
holds, for all $\epsilon\in (0,1]$, where $u_\epsilon^\star$ is in~\eqref{eq:optimal-control}. Thus, from~\eqref{eq: finite_energy_bound} we have that 
\begin{multline*}
\int_Q \frac12\Big\|\frac{d\mathbf m^*_\epsilon(t,x)}{d\boldsymbol\mu^*_\epsilon(t,x)}\Big\|^2 d\boldsymbol\mu^*_\epsilon\\\leq \int_{\R^d}\rho_f(y)\log(\rho_f(y))dy-\int_{\R^d\times\R^d}\epsilon\log(p_{t_f,\epsilon}(x,y))\rho_0(x)\rho_f(y)dxdy,    
\end{multline*}
 for all $\epsilon\in (0,1]$. Since $\rho_0,\rho_{f}$ are compactly supported, from~\cite[Theorem~19]{barilari2012small}, we have that $\{\epsilon\log(p_{t_f,\epsilon}(x,y))\}_{\epsilon\in(0,1]}$ is uniformly bounded. Hence, we get the uniform bound
 \begin{equation}
    \label{eq:uniform_energy_explicit}
    \int_Q \frac12\Big\|\frac{d\mathbf m^*_\epsilon(t,x)}{d\boldsymbol\mu^*_\epsilon(t,x)}\Big\|^2 d\boldsymbol\mu^*_\epsilon=\int_Q  \frac{\|m_{\epsilon}^\star(t,x)\|^2}{2\rho_\epsilon^*(t,x)} dx\, dt \leq C,
\end{equation}
where $C$ is independent of $\epsilon$. Hence, the Cauchy--Schwarz inequality, for any Borel set $B \subset Q$ yields:
\begin{equation}
    |\mathbf m^*_\epsilon|(B) = \int_B \Big\|\frac{d\mathbf m^*_\epsilon(t,x)}{d\boldsymbol\mu^*_\epsilon(t,x)}\Big\| d\boldsymbol\mu^*_\epsilon \leq \left( 2C \right)^{1/2} \boldsymbol\mu^*_\epsilon(B)^{1/2}.
\end{equation}
Thus tightness of $\{\boldsymbol\mu^*_\epsilon\}_\epsilon$ implies tightness of $\{\mathbf m^*_\epsilon\}_\epsilon$. Since $Q=[0,t_f]\times\R^d$ and $\boldsymbol\mu^*_\epsilon(dt,dx)=\mu^*_{\epsilon,t}(dx)dt$, it is enough to show tightness of $\mu^*_{\epsilon,t}$. To this end, let $\mathrm{P}^{u_{\epsilon}^*}$ be the probability distribution of the optimal process $ X^{u^*_{\epsilon}}\in \Omega$ in~\eqref{controlled_dynamics}, where $u^*_\epsilon$ is in~\eqref{eq:optimal-control}.
Using the variational formula for relative entropy~\cite[\S6.6]{DemboZeitouni1998},\cite[\S 1.4]{DupuisEllis1997},
\[
D_{\mathrm{KL}}(\mathrm{P}^{u_{\epsilon}^*} \| R_{\epsilon}) = \sup_{f} \left\{ \int f d\mathrm{P}^{u_{\epsilon}^*} - \log \int e^{f} dR_{\epsilon} \right\},
\]
we get that 
\[
\mathrm{P}^{u_{\epsilon}^*}(K_L^c)\ \le\frac{ \KL(\mathrm{P}^{u_{\epsilon}^*}\|\mathrm{R}^{\epsilon,\mu_0})+\log 2}{\log\big(1/\mathrm{R}^{\epsilon,\mu_0}(K_L^c)\big)},\qquad\text{for all $\epsilon>0$}
\]
%{\color{blue}
where $\{K_L\}_{L\in\mathbb{N}}\subset\Omega$ are compact sets. Following from~\cite[Theorem~4.3]{BoueDupuis1998}, since $\mathrm R^\epsilon $ is exponential tight (see for example~\cite[pp~8]{DemboZeitouni1998} for definition of exponential tightness), together with~\eqref{eq:uniform_energy_explicit}, we have that 
\[
\mathrm{P}^{u_{\epsilon}^*}(K_L^c)\ \le\frac{C+\log 2}{L},\qquad\text{for all $\epsilon>0$}
\]
and hence 
\[
\sup_{\epsilon>0} \mathrm{P}^{u_{\epsilon}^*}(K_L^c)\ \le\ \frac{C+\log 2}{L}\xrightarrow[L\downarrow \infty]{}0.
\]
Hence, by Prokhorov Theorem~\cite[Theorem~5.1]{Billingsley1995}, we get that $\{\mathrm{P}^{u_{\epsilon}^*}\}_{\epsilon>0}$ is tight, which concludes that $\mu^*_{\epsilon,t}=\mathrm{ev}_{t\#}\mathrm{P}^{u_{\epsilon}^*}$ is tight. Therefore, by Prokhorov's theorem, $(\boldsymbol{\mu}_\epsilon^*, \mathbf{m}_\epsilon^*)$ is precompact. Thus, there exists a subsequence (not relabeled) $(\boldsymbol{\mu}_\epsilon^*, \mathbf{m}_\epsilon^*) \in \mathcal{M}_+(\overline{Q}) \times \mathcal{M}(\overline{Q}; \mathbb{R}^m)$ satisfying % the Fokker-Planck system in the weak sense as space-time measures:
\begin{multline}\label{eq:FP-weak}
\int_Q\!\Big(
-\partial_t\phi(t,x)\boldsymbol\mu_\epsilon(dt,dx)
- \nabla_x\phi(t,x)\cdot \big(b_\epsilon(x)\boldsymbol\mu_\epsilon(dt,dx) + g(t,x)\,\mathbf m_\epsilon(dt,dx)\big) \\
- \tfrac{\epsilon}{2}\sum_{i,j}\partial^2_{x_ix_j}\phi(t,x) G_{ij}(t,x)\,\boldsymbol\mu_\epsilon(dt,dx)
\Big) \\+\!\int_{\mathbb R^d}\!\phi(0,x)\,\mu_0(dx)-\int_{\mathbb R^d}\!\phi(t_f,x)\,\mu_f(dx)=0    
\end{multline}
for all $\phi\in C_c^\infty((0,t_f)\times\mathbb R^d)$ and
\begin{equation}
    \label{eq:weak_convergence}
    (\boldsymbol{\mu}_\epsilon^*, \mathbf{m}_\epsilon^*) \rightharpoonup (\boldsymbol{\mu}^*, \mathbf{m}^*)\qquad \text{ in }\qquad \mathcal{M}_+(\overline{Q})\times\mathcal{M}(\overline{Q}; \mathbb{R}^m).
\end{equation}
Pass the limit in the integral yields that $(\boldsymbol{\mu}^*, \mathbf{m}^*)$ satisfies~\eqref{eq:CE_measure_unref}. Since $\mathcal{J}$ is strictly convex and lower semicontinuous with respect to weak$^*$ convergence for every $\epsilon\geq 0$, we have that
\begin{equation}
    \label{eq:lsc_cost}
    \liminf_{\epsilon \to 0}\mathcal{J}(\boldsymbol{\mu}_\epsilon^*, \mathbf{m}_\epsilon^*) \geq \mathcal{J}(\boldsymbol{\mu}^*, \mathbf{m}^*).
\end{equation}

%{\color{red}
%\begin{itemize}
%\item What is the standard mollifier $\eta_{\varepsilon}$? And why are we mollifying? The $\varepsilon>0$ problem is well-posed for compactly supported measures from what I understand.
%\item Just a remark: the measures $\tilde{\mu}^{\varepsilon}_{t_i}$ don't have compact support if $\eta_{\varepsilon}$ does not have compact support. Following like 435-437, we seemed to need compactness of supports for the $\varepsilon-$ problem.
%\item Lower semicontinuity of the cost function on the paths and coercivity is not obvious. Standard argument for coercivity in fact lost because of the degeneracy of the metric. 
%\item I think I am not following what the splitting of time intervals below is achieving. It feels like if we take $t_1 = t_f$, the argument still works. We could potentially replace $\tilde{\mu}, \tilde{m}$ by an optimal solution to the limit problem, and then apply large deviations result?

%\item $[39]$ is working in the Riemannian setting. Though, I think the conclusion is true.
%\item I don't follow why 4.14 holds true below. Why is the sum of KL divergences for the $\epsilon>0$ problem less than the sum of the OT costs? Are we using some property of KL divergence here?
%\item Line 518:"Following from [37, Proposition 2.5, Corollary 2.6, Theorem 2.7]", This feels like citing too many results at once. Is the final one not enough?
%\end{itemize}
%}
To conclude that $(\boldsymbol{\mu}^*, \mathbf{m}^*)$ is a minimizer, we prove the corresponding lower limit superior bound for~\eqref{eq:lsc_cost}. To this end, we consider two functionals
\[
\mathcal G_{\epsilon}(\mathrm{P}):= \epsilon D_{\mathrm{KL}}(\mathrm{P} \| \mathrm{R}^{\epsilon,\mu_0}) + \iota_{\{(\mathrm{ev}_0)_\# \mathrm{P}=\mu_0,\; (\mathrm{ev}_{t_f})_\# \mathrm{P}=\mu_f\}}(\mathrm{P})
\]
where $\mathrm{R}^{\epsilon,\mu_0}$ is the law of~\eqref{uncontrolled_dynamic_constraint} and
\[
\mathcal G_{0}(\mathrm{P}):=\int_{\Omega} \mathrm{C}(\omega) \, d\mathrm{P} + \iota_{\{(\mathrm{ev}_0)_\# \mathrm{P}=\mu_0,\; (\mathrm{ev}_{t_f})_\# \mathrm{P}=\mu_f\}}(\mathrm{P}),
\]
where
\begin{equation}\label{eq: rate_function}
\mathrm{C}(\omega):=
\begin{cases}
\displaystyle
\frac12\int_{0}^{t_{f}}
\big\langle \dot\omega_t,\,
G(\omega_t)^\dagger \dot\omega_t\big\rangle\,dt,
&
\omega \in \mathrm{AC}([0,t_f]; \mathbb{R}^d), \; \dot{\omega}_t \in \mathcal{D}_{\omega_t} \text{ a.e.}, \\[10pt]
+\infty, & \text{otherwise}.
\end{cases}
\end{equation}

We claim that 
\begin{equation}\label{eq: Gamma_convergence}
\Gamma\text{-}\lim_{\epsilon \rightarrow 0} \mathcal G_{\epsilon}= \mathcal G_0   \qquad\text{on}\qquad\mathcal P(\Omega).
\end{equation}
Once this is established, the recovery part of the definition of $\Gamma$-convergence~\cite{dal2012introduction} will provide us with the lower limit superior. To prove~\eqref{eq: Gamma_convergence}, we use large deviation principle (LDP)~\cite{DemboZeitouni1998,DupuisEllis1997} for  the law of~\eqref{uncontrolled_dynamic_constraint} initialized at $X_{0}=x$, which we denote as $\mathrm{R}^{\epsilon,x}\in\mathcal P(\Omega)$. To this end, since $b_{\epsilon} :=\frac\epsilon2\sum_{i=1}^{m}\nabla_{g_i}g_i\rightarrow 0$ as $\epsilon\rightarrow0$  uniformly on $\mathbb{R}^d$ and under Assumption~\ref{ass:wellposed} we have $b_{\epsilon} :=\frac\epsilon2\sum_{i=1}^{m}\nabla_{g_i}g_i$ satisfies linear growth,  following from~\cite[Theorem~3.1]{ChiariniFischer2014}, we have that $\mathrm{R}^{\epsilon,x}$ satisfies large deviation principle with good rate functional
\[
I_x(\omega) = \inf_{\substack{ u \in L^2([0,t_f]; \mathbb{R}^m) \\ \dot{\omega}_t = g(\omega_t) u_t \text{ a.e.} \\ \omega_0 = x }} \frac12 \int_0^{t_f} \|u_t\|^2 \, dt=\frac12 \int_0^{t_f} \langle \dot{\omega}_t, G(\omega_t)^\dagger \dot{\omega}_t \rangle \, dt=  C_x(\omega),
\] 
where $C_x(\omega):= \mathrm{C}(\omega) + \iota_{\{\omega_0 = x\}}(\omega)$ and $\mathrm{C}$ is defined in~\eqref{eq: rate_function}. Therefore, following from~\cite[Proposition~2.5 and Proof of Theorem~6.1]{Leonard2012} we conclude that~\eqref{eq: Gamma_convergence} holds. 

Let $(\tilde{\boldsymbol{\mu}}, \tilde{\mathbf{m}})$ be any admissible pair satisfying \eqref{eq:CE_measure_unref} with $\mathcal{J}(\tilde{\boldsymbol{\mu}}, \tilde{\mathbf{m}}) < \infty$. Then, following from~\cite[Theorem~8.2.1]{ambrosio2005gradient}, there exists a probability measure $\mathrm{\tilde P} \in \mathcal{P}(\Omega)$  such that 
\begin{equation}\label{eq: path-space_representation}
\tilde{\boldsymbol{\mu}}_t= (\mathrm{ev}_t)_\# \mathrm{\tilde P}.
\end{equation}
Moreover, by setting $\tilde{\mathbf{m}} = \tilde{u} \tilde{\boldsymbol{\mu}}$, the path measure $\mathrm{\tilde P}$ is concentrated on absolutely continuous $\omega$ satisfying $\dot{\omega}_t = g(\omega_t) \tilde{u}(t,\omega_t)\in \mathcal{D}_{\omega_t}$ in~\eqref{eq: distribution} a.e $t\in [0,t_f]$  and
\begin{equation}\label{eq: measures_to path_measures}
\mathcal{J}(\tilde{\boldsymbol{\mu}}, \tilde{\mathbf{m}}) = \int_\Omega \int_0^{t_f} \frac12 \|\tilde{u}(t,\omega_t)\|^2 dt d\mathrm{\tilde P}<\infty.  
\end{equation}
Since $(\mathrm{ev}_0,\mathrm{ev}_{t_f})_\#\mathrm{\tilde P}=(\mu_0,\mu_f)$  and $\mathcal G_0(\mathrm{\tilde P})= \int_{\Omega} \mathrm{C}(\omega) \, d\mathrm{\tilde P}=\mathcal{J}(\tilde{\boldsymbol{\mu}}, \tilde{\mathbf{m}})<\infty$, the recovery part of~\eqref{eq: Gamma_convergence} implies that there exists a sequence 
$\mathrm{P}^\epsilon\rightharpoonup\mathrm{\tilde P}$ with 
$(\mathrm{ev}_0,\mathrm{ev}_{t_f})_\#\mathrm{P}^\epsilon=(\mu_0,\mu_f)$ such that
\begin{equation}
    \label{eq:recovery_path_space}
    \limsup_{\epsilon \to 0} \epsilon D_{\mathrm{KL}}(\mathrm{P}^\epsilon \| \mathrm{R}^\epsilon) \leq \int_{\Omega} \mathrm{C}(\omega) \, d\mathrm{\tilde P}=\mathcal{J}(\tilde{\boldsymbol{\mu}}, \tilde{\mathbf{m}})<\infty.
\end{equation}
This proves the claim. 

We are finally in a position to prove~\eqref{eq:aux1}. Since $\mathrm{P}^\epsilon \ll \mathrm{R}^\epsilon$, by Girsanov's theorem~\cite[Chapter~8.6]{oksendal2003stochastic} we have that
\[
\left.\frac{d\mathrm{R}^\epsilon}{d\mathrm{P}^\epsilon}\right|_{\mathcal{F}_t} = \exp\left( \int_0^t \frac{1}{\sqrt{\epsilon}} \tilde{u}_\epsilon^\top dW_s - \frac12 \int_0^t \frac{1}{\epsilon} \|\tilde{u}_\epsilon\|^2 ds \right)
\]
for some adapted process $\tilde{u}_\epsilon$. This implies that under $\mathrm{P}^\epsilon$, the canonical process satisfies:
\[
dx_t = \bigl( b_\epsilon(x_t) + g(x_t) \tilde{u}_\epsilon(t, x_t) \bigr) dt + \sqrt{\epsilon} \, g(x_t) dW_t^{\mathrm{P}^\epsilon}
\]

%We state here that, under Assumption~\ref{ass:Controllability condition} and~\ref{ass:wellposed}, one can construct this sequence using similar techniques in the Euclidean setting in~\cite[Proposition 2.5]{Leonard2012}. The key difference here is that we rely on the sub-Riemannian metric which exists under our controllability assumption. {\color{red} For the purpose of space we postpone this derivation to another avenue.}\margin{So, I don't like this whole paragraph as part of the proof. Do we need to say this here? \kecmt{ke:I agree.}} 
%Define $(\tilde{\rho}_\epsilon, \tilde{m}_\epsilon)$ from $\mathrm{P}^\epsilon$ via the Eulerian disintegration:
From the sequence $\{\mathrm{P}^\epsilon\}_\epsilon$, we construct the sequence of pairs %\margin{pairs} 
$(\tilde{\boldsymbol{\mu}}_\epsilon, \tilde{\mathbf{m}}_\epsilon)$ via the pushforward mappings% using the continuous mappings
\begin{equation}
    \label{eq:eulerian_from_path}
    \tilde{\boldsymbol{\mu}}_{\epsilon,t} := (\mathrm{ev}_t)_\# \mathrm{P}^\epsilon, \qquad \int_Q \psi \cdot d\tilde{\mathbf{m}}_\epsilon := \mathbb{E}^{\mathrm{P}^\epsilon}\left[ \int_0^{t_f} \psi(t, x_t) \cdot \tilde{u}_\epsilon(t, x_t) dt \right],
\end{equation}
%\margin{A bit confused here because I thought that we show $(x_t)_\# \mathrm{P}^\epsilon$ by $ \tilde{\boldsymbol{\mu}}_{\epsilon,t}$. Maybe I am missing something here}
where $\psi\in C_c(Q;\R^d)$. Then $(\tilde{\boldsymbol{\mu}}_\epsilon, \tilde{\mathbf{m}}_\epsilon)$ satisfies~\eqref{eq: fb_measure_constraint}-\eqref{eq:endpoint_measure_constraints} and %\margin{are these correct eqref ?} and by construction:
\[
\mathcal{J}(\tilde{\boldsymbol{\mu}}_\epsilon, \tilde{\mathbf{m}}_\epsilon) = \epsilon D_{\mathrm{KL}}(\mathrm{P}^\epsilon \| \mathrm{R}^\epsilon).
\]
By optimality, we have that $\mathcal{J}(\boldsymbol{\mu}_\epsilon^*, \mathbf{m}_\epsilon^*)\leq \mathcal{J}(\tilde{\boldsymbol{\mu}}_\epsilon, \tilde{\mathbf{m}}_\epsilon)$ and hence
\begin{equation}
    \mathcal{J}(\boldsymbol{\mu}^*, \mathbf{m}^*) \leq \liminf_{\epsilon \to 0} \mathcal{J}(\boldsymbol{\mu}_\epsilon^*, \mathbf{m}_\epsilon^*) \leq \limsup_{\epsilon \to 0} \mathcal{J}(\tilde{\boldsymbol{\mu}}_\epsilon, \tilde{\mathbf{m}}_\epsilon) \leq \mathcal{J}(\tilde{\boldsymbol{\mu}}, \tilde{\mathbf{m}}).
\end{equation}
\end{proof}

\section{Numerical Method and Examples}\label{sec: Numerical Method}
In this section, we develop numerical methods for solving the optimality conditions~\eqref{eq:Schr-fwd}-\eqref{eq:Schr-bc}. %{\color{blue}Unlike the purely hyperbolic transport equations in distributional optimal transport, our parabolic formulation enables robust forward-backward numerical schemes.} 
Using the boundary condition in~\eqref{eq:Schr-bc} we can define the operator $\map{T_{\epsilon}}{L^2(\R^d)}{L^2(\R^d)}$ as
\begin{equation}\label{eq: contraction operator}
(T_{\epsilon}\varphi_f)(y)
:=
\frac{\rho_f(y)}{\displaystyle\int_{\mathbb R^d} p_{t_f,\epsilon}(x,y)
\left[ \frac{\rho_0(x)}{(Q_{t_f}\varphi_f)(x)} \right] dx },
\end{equation}
%\margin{Where is $Q_{t_f}$ defined?}
where $Q_{t_f}$ is defined in~\eqref{eq: action_semigroup}. Since $\rho_0,\rho_f$ are positive density, compactly supported (or having exponential tail) on $\R^d$ and the transition kernel $p_{t,\epsilon}(x,y)$ is 
strictly positive under the bracketing  Assumption in~\ref{ass:Controllability condition} (or the H\"ormander condition), following from~\cite{Chen2016Hilbert}, $T_{\epsilon}$ in~\eqref{eq: contraction operator} is a contraction in the Hilbert metric on the cone of positive functions; hence the iteration 
$\varphi_f^{k+1} = T_{\epsilon}(\varphi_f^{k})$ converges linearly to the unique
fixed point $\varphi_f$, where 
\[
T_{\epsilon}(\varphi_f)= \varphi_f.
\]
Once the fixed point $\varphi_f$ is obtained, the full space-time solution for~\eqref{eq:Schr-fwd}-\eqref{eq:Schr-bc} is recovered by
\[
\varphi(t,x) = (Q_{t_f-t} \varphi_f)(x), \qquad\text{and}\qquad
\widehat{\varphi}(t,x) = (P_t\big( \rho_0 / (Q_{t_f}\varphi_f) \big))(x),
\]
where $P_t$ and $Q_{t_f}$ is defined in~\eqref{eq: action_semigroup}. We state here that, for \(\epsilon>0\), the transition density \(p_{t,\epsilon}(x,y)\)
is the fundamental solution of the forward Kolmogorov equation in the $y$-variable:
\[
\partial_t p_{t,\epsilon}(x,y) - (\mathcal L_{\epsilon,y}^* p_{t,\epsilon})(x,y)=0,\qquad t>0,
\qquad
p_{0,\epsilon}(x,y)=\delta_x(y),
\]
where \(\mathcal L_{\epsilon,y}^*\) denotes \(\mathcal L_\epsilon^*\) acting on the \(y\)-variable.
%Equivalently, \(p_{t,\epsilon}\) satisfies the backward Kolmogorov equation in the \(x\)-variable:
%\[
%\partial_t p_{t,\epsilon}(x,y) + (\mathcal L_{\epsilon,x} p_{t,\epsilon})(x,y)=0,\qquad t>0\qquad
%p_{0,\epsilon}(x,y)=\delta_y(x).
%\]
\begin{figure}[t]
    \centering
    
    \begin{subfigure}[b]{0.45\textwidth}
        \centering
        \includegraphics[width=0.5\textwidth]{rho_0.png}
        \caption{Initial Gaussian $\rho_0$}
        \label{fig:rho0}
    \end{subfigure}
    \hfill
    \begin{subfigure}[b]{0.45\textwidth}
        \centering
        \includegraphics[width=0.5\textwidth]{rho_f.png}
        \caption{Final ring $\rho_f$}
        \label{fig:rhof}
    \end{subfigure}
    \caption{The above is the plot of the  isosurfaces of the initial and final densities defined in~\eqref{eq: initial density} and~\eqref{eq: final density}, respectively.}
    \label{fig:boundary_densities}
\end{figure}
%\margin{where is tol defined/initialized below?}
\begin{algorithm}
{\footnotesize
\caption{Static Sinkhorn for Degenerate Schr\"odinger Bridge}
\begin{enumerate}
\item \textbf{Input:} $\rho_0,\rho_f$, tolerance level $\text{tol}$, $\{\epsilon_0 > \cdots > \epsilon_K\approx 0\}$, 
      transition kernel $p_{t,\epsilon_k}(x,y)$ % for each $\epsilon$

\item \textbf{For each $\epsilon_k$ in $\{\epsilon_0, \epsilon_1, \dots, \epsilon_K \approx 0\}$:}
\begin{enumerate}
\item If $k > 0$: initialize $\varphi_f^{\epsilon_k,(0)} = \varphi_f^{\epsilon_{k-1},*}$
\item \textbf{Sinkhorn Iteration until convergence:}
For $m \in \mathbb{N}_{\geq 0}$:
\begin{enumerate}
\item \textbf{Compute $\widehat{\varphi}_0^{\epsilon_k,m}$:}
\[
\widehat{\varphi}_0^{\epsilon_k,m}(x) = \frac{\rho_0(x)}{\displaystyle\int p_{t_f,\epsilon_k}(x,y) \varphi_f^{\epsilon_k,m}(y)   dy}
\]

\item \textbf{Update $\varphi_f^{\epsilon_k,m+1}$:}
\[
\varphi_f^{\epsilon_k,m+1}(y) = \frac{\rho_f(y)}{\displaystyle\int p_{t_f,\epsilon_k}(x,y) \widehat{\varphi}_0^{\epsilon_k,m}(x)   dx}
\]

\item \textbf{Check convergence:} $\|\varphi_f^{\epsilon_k,m+1} - \varphi_f^{\epsilon_k,m}\| < \text{tol}$
\end{enumerate}

\item \textbf{Store fixed point:} $\varphi_f^{\epsilon_k,*} = \varphi_f^{\epsilon_k,m}$

\item \textbf{Reconstruct full solution via single PDE solves:}
\begin{enumerate}
\item \textbf{Store $\varphi_{\epsilon_k}^*(t,x)$ from backward equation:}
\[
\partial_t\varphi + \mathcal{L}_{\epsilon_k}\varphi = 0,\quad \varphi(t_f,\cdot) = \varphi_f^{\epsilon_k,*}
\]


\item \textbf{Compute initial $\widehat{\varphi}$ from $\widehat{\varphi}_{\epsilon_k}^*(0,x) = \rho_0(x) / \varphi_{\epsilon_k}^*(0,x)$ and store $\widehat{\varphi}_{\epsilon_k}^*(t,x)$ from the Forward equation:}
\[
\partial_t\widehat{\varphi} - \mathcal{L}_{\epsilon_k}^*\widehat{\varphi} = 0,\quad 
\widehat{\varphi}(0,\cdot) = \widehat{\varphi}_{\epsilon_k}^*(0,\cdot)
\]

\end{enumerate}

\item \textbf{Compute optimal quantities $\{\rho_{\epsilon_k}^*, u_{\epsilon_k}^*\}_{k=0}^K$:}
\begin{align*}
\rho_{\epsilon_k}^*(t,x) = \widehat{\varphi}_{\epsilon_k}^*(t,x) \varphi_{\epsilon_k}^*(t,x)\quad\text{and}\quad
u_{\epsilon_k}^*(t,x) = \epsilon_k   g(x)^\top \nabla \log \varphi_{\epsilon_k}^*(t,x)
\end{align*}
\end{enumerate}
\item \textbf{Simulate controlled dynamics for verification:}
For each $\epsilon_k$, simulate:
\[
dX_t^u = g(X_t^u) u_{\epsilon_k}^*(t, X_t^u) dt + \sqrt{\epsilon_k} g(X_t^u) dW_t,
\quad x_0^u \sim \mu_0
\]
\end{enumerate}
}
\end{algorithm}

\begin{figure}[t]
    \centering
    \vspace{-0.3cm}
    % Epsilon = 1
    \includegraphics[width=0.9\textwidth]{rho_t_e1.png}
    
    \vspace{-0.3cm}
    
  %  % Epsilon = 0.5
  %  \includegraphics[width=0.9\textwidth]{rho(t)_e0.5.png}
    
  %  \vspace{-0.3cm}
  %  
  %  % Epsilon = 0.1
  %  \includegraphics[width=0.9\textwidth]{rho(t)_e0.1.png}
    
 %   \vspace{-0.3cm}
    
    % Epsilon = 0.01
    \includegraphics[width=0.9\textwidth]{rho_t_e0.01.png}
    
    \caption{The plots above are the isosurface plots of the optimal density $\rho(t)$ associated to Example~\ref{ex: unified_example}, with $\alpha=\frac{1}{4}$. The aim is to demonstrate a smooth and continuous morphing from the initial Gaussian blob in~\eqref{eq: initial density} to the final ring-shaped distribution in~\eqref{eq: final density}, for any $\epsilon>0$. Hence, as $\epsilon \to 0$, the Schr\"{o}dinger bridge approaches deterministic optimal transport.}
    \label{fig:all_eps_morphing}
\end{figure}
\begin{figure}[t]
  \centering
  \includegraphics[width=\textwidth]{Continuous_SDE_bridge_paths.png}
  \caption{This plot is the corresponding Heisenberg bridge path associated to Example~\ref{ex: unified_example}. Figure~\ref{fig:all_eps_morphing} demonstrates the morphing from the initial to final densities. This is the
sample trajectories of the controlled diffusion used to realize a Schr\"odinger bridge on $3D$
domain for noise intensities
$\epsilon\in\{1,0.5,0.1,0.01\}$.
Black curves show $n=50$ (for better visualization) sample paths $t\mapsto(x_t,y_t,z_t)$ in~\eqref{controlled_dynamics} where~\eqref{eq: diff_coefficient}; marked points indicate the initial and
terminal positions.
The bridge dynamics is obtained by changing the drift of the uncontrolled horizontally constrained Heisenberg-type reference
diffusion in Fig.~\ref{fig:diffusion_comparison_three}-(a) so that its time-marginals interpolate
between prescribed endpoint laws.
As $\epsilon$ decreases, the ensemble concentrates
and trajectories exhibit constrained curvature indicating the effect of the presence of horizontal constraint.}
\label{fig:heisenberg-sb-paths}
\end{figure}
%\margin{We have a subsection of Examples and they example environment. I would remove/change this subsection title}
%\subsection{Examples}
We state here that the minimum dimension for our setup to be non-trivial is $3$. The reason is that in this setting, one can genuinely have non-holonomic constraints with underactuated systems and the  Schr\"odinger bridge problem, interpreted as the optimal change in drift
becomes non-trivial. Throughout, we fix the initial density to be the standard Gaussian
 \begin{equation}\label{eq: initial density}
 \rho_0(x,y,z)
=
\frac{1}{(2\pi)^{3/2} \sigma_x\sigma_y\sigma_z} 
\exp\!\left(
-\frac12\left[
\frac{(x-m_x)^2}{\sigma_x^2}
+\frac{(y-m_y)^2}{\sigma_y^2}
+\frac{(z-m_z)^2}{\sigma_z^2}
\right]\right),     
 \end{equation}
where $(m_x,m_y,m_z)=(0,0,0)$, $\sigma_x=\sigma_y=\sigma_z=0.55$, and the desired density to be ``ring-like'' (see, Figure~\ref{fig:boundary_densities}).%\margin{should you refer to the figure so that people know what ring like means? }:
\begin{equation}\label{eq: final density} 
\rho_f(x,y,z)=\frac{1}{Z_f} \tilde\rho_f(x,y,z),
\qquad
Z_f=\iiint_{\mathbb{R}^3}\tilde\rho_f(x,y,z) dx dy dz,    
\end{equation}
where 
\[
\tilde\rho_f(x,y,z)
=
\exp\!\left(-\frac{(r-R_0)^2}{2s_R^2}\right) 
\exp\!\left(-\frac{(z-m_z)^2}{2\sigma_z^2}\right),
\qquad r=\sqrt{x^2+y^2}.
\]
Here $r=\sqrt{x^2+y^2}$,  $R_0=2.5$, $s_R=0.45$, $m_z=0$, and $\sigma_z=0.70$ (see Figure~\ref{fig:boundary_densities}).
We illustrate our theoretical results above by the following examples.

%demonstrate this by considering three fundamental examples of optimal transport over non-holonomic constraints: the Brockett integrator, the Heisenberg system, and the Unicycle model which illustrate our theoretical results above.
%[Family of Sub-Riemannian Schr\"odinger Bridges]
%\margin{Using \em{} change the example environment from italic to standard, everywhere}
\begin{example}\label{ex: unified_example}
\em{Given initial and final densities $\rho_0,\rho_f$ in~\eqref{eq: initial density} and~\eqref{eq: final density}, we consider the optimal transport from an ensemble of states   $\begin{pmatrix}x_0 \\ y_0 \\ z_0\end{pmatrix} \sim \rho_0$, where each particle's motion is constrained by the following controlled dynamics:
\begin{equation}\label{eq: unified_dynamics}
d\begin{pmatrix}x_t \\ y_t \\ z_t\end{pmatrix} = \begin{pmatrix}1 & 0 \\ 0 & 1 \\ -2\alpha y_t & 2\alpha x_t\end{pmatrix}\begin{pmatrix}u^1_t \\ u^2_t\end{pmatrix}dt,
\end{equation}
where $\alpha>0$, so that at time $t=t_f$ we have that $\begin{pmatrix}x_{t_f} \\ y_{t_f} \\ z_{t_f}\end{pmatrix} \sim \rho_f$. Here optimality is measured in~\eqref{eq: expectation_total_control_effort}. The vector fields associated with the control directions are:
\begin{equation}\label{eq: unified_vector_fields}
g_1 = \partial_{x} - 2\alpha y \partial_z, \quad g_2 = \partial_{y} + 2\alpha x \partial_z.
\end{equation}

To formulate the associated Schr\"{o}dinger bridge problem, we introduce a small noise aligned with the control directions, leading to the stochastic process~\eqref{controlled_dynamics}
where 
\begin{equation}\label{eq: diff_coefficient}
g(x,y,z) = \begin{pmatrix}1 & 0 \\ 0 & 1 \\ -2\alpha y & 2\alpha x\end{pmatrix}\qquad \text{and}\qquad\text{ $W_t$ is a Brownian motion in $\mathbb{R}^2$}.    
\end{equation}
 The diffusion tensor is given by:
\begin{equation}\label{eq: unified_diffusion_tensor}
G = gg^\top = \begin{pmatrix}
1 & 0 & -2\alpha y \\
0 & 1 & 2\alpha x \\
-2\alpha y & 2\alpha x & 4\alpha^2(x^2+y^2)
\end{pmatrix},
\end{equation}
which has $\mathrm{rank}(G) = 2 < 3$ and is thus degenerate. However, the Lie bracket of the vector fields yields:
\[
[g_1, g_2] = 4\alpha \partial_z,
\]
which is everywhere linearly independent of $g_1$ and $g_2$. Therefore, $$\mathrm{span}\{g_1, g_2, [g_1, g_2]\} = \mathbb{R}^3,$$ satisfying the H\"ormander condition (Assumption~\ref{ass:Controllability condition}).

Let us define the scaled vector fields for the noise:
\begin{equation*}
V_1 = \sqrt{\epsilon}(\partial_x - 2\alpha y\partial_z), \quad\text{and}\quad 
V_2 = \sqrt{\epsilon}(\partial_y + 2\alpha x\partial_z).
\end{equation*}
The drift vector field $V_0$ is given by:
\begin{align*}
V_0 &= -\frac{\epsilon}{2}[\nabla_{g_1}g_1 + \nabla_{g_2}g_2]\cdot\nabla \\
&= -\frac{\epsilon}{2}[(0,0,0)^\top + (0,0,0)^\top]\cdot\nabla = 0.
\end{align*}
Thus, the infinitesimal generator for the process~\eqref{controlled_dynamics} satisfies:
\begin{equation}\label{eq: unified_infinitesimal_generator}
\mathcal{L}_\epsilon = \frac12(V_1^2 + V_2^2)
\end{equation}
and we have:
\[
\mathrm{Lie}(V_1,V_2) = \mathrm{span}\{g_1, g_2, [g_1,g_2]\} = \mathbb{R}^3.
\]
This implies that the parabolic operators
\[
\partial_t + \mathcal{L}_\epsilon
\quad\text{and}\quad
\partial_t - \mathcal{L}_\epsilon^*
\]
with $\mathcal{L}_\epsilon$ as in~\eqref{eq: unified_infinitesimal_generator} are hypoelliptic.
In particular, there exists a smooth and positive transition density $p_{t,\epsilon}$ that solves 
\begin{equation}\label{eq: forw_equ}
(\partial_t - \mathcal{L}_\epsilon^*)\rho = 0.
\end{equation}
To obtain an explicit solution for the transition density function $p_{t,\epsilon}$, we rely on the group structure associated to the vector fields. Let $\mathbb{H}:=(\mathbb{R}^3,\cdot)$ with group multiplication:
\[
(x,y,z) \cdot (x',y',z') = (x + x', y + y', z + z' + 2\alpha(x y' - y x'))
\]
for all $(x,y,z),(x',y',z') \in \mathbb{R}^3$. Then, one can check that the  vector fields in~\eqref{eq: unified_vector_fields} are left-invariant on $\mathbb{H}$. Thus, following the same calculations in~\cite{barilari2012small,bonfiglioli2007stratified}, we get that
\begin{equation}
p_{t,\epsilon}(\vec{0}, \vec{q}) = \frac{1}{(2\pi \epsilon t)^2} \int_{\mathbb{R}} \frac{2\alpha\lambda}{\sinh(2\alpha\lambda \epsilon t)} 
\exp\left(-\frac{\alpha\lambda(x^2 + y^2)}{\tanh(2\alpha\lambda \epsilon t)} + 2i\alpha\lambda z\right) d\lambda
\end{equation}
is the transition density function from $\vec{0} = (0,0,0)$ at time $0$ to $\vec{q}_t = (x,y,z)$ at time $t$ that solves~\eqref{eq: forw_equ} with initial condition $p_{0,\epsilon}(\vec{0}, \vec{q}) = \delta_{\vec{0}}(\vec{q})$.
 Therefore, using the group structure, we translate the initial identity to any arbitrary point. That is, for any $\vec{q}=(x,y,z)$ and $\vec{q}_0=(x^\prime,y^\prime,z^\prime)$ we have that 
 \[
p_{t,\epsilon}(\vec{q}_0, \vec{q})=p_{t,\epsilon}(\vec{0}, \vec{q}_0^{-1}\cdot\vec{q}), 
\]
where $(\vec{q}_0)^{-1} =(-x^\prime, -y^\prime, -z^\prime)$. Therefore, the transition density function from $\vec{q}_0=(x^\prime,y^\prime,z^\prime)$ at time $t=0$ to $\vec{q}=(x,y,z)$ at time $t$ is
\begin{multline}\label{eq: unified_heat_kernel}
p_{t,\epsilon}(\vec{q}_0, \vec{q}) = \frac{1}{(2\pi\epsilon t)^2} \int_{\mathbb{R}} \frac{2\alpha\lambda}{\sinh(2\alpha\lambda \epsilon t)} 
\exp\bigg(-\frac{\alpha\lambda((x-x^\prime)^2 + (y-y^\prime)^2)}{\tanh(2\alpha\lambda \epsilon t)} \\
+ 2i\alpha\lambda (z - z^\prime - 2\alpha(x^\prime y - y^\prime x))\bigg) d\lambda.
\end{multline}

Let $(\phi_0,\phi_f)$ satisfy the  Schr\"odinger system 
\begin{equation}\label{eq: SS_BI}
\phi_0(\vec{q}_0)(\mathscr P_{0,t_f} \phi_f)(\vec{q}_0)= \rho_0(\vec{q}_0), \quad\text{and}\quad
(\mathscr P_{0,t_f}^* \phi_0)(\vec{q}_1)\phi_f(\vec{q}_1)= \rho_f(\vec{q}_1),    
\end{equation}
where
\begin{align*}
(\mathscr P_{0,t_f} \phi_f)(\vec{q}_0) =& \int \phi_f(\vec{q}_1) p_{t_f,\epsilon}(\vec{q}_0, \vec{q}_1) d\vec{q}_1, \quad \text{and}\\ 
(\mathscr P_{0,t_f}^* \phi_0)(\vec{q}_1)=& \int \phi_0(\vec{q}_0) p_{t_f,\epsilon}(\vec{q}_0, \vec{q}_1) d\vec{q}_0.    
\end{align*}
Then, following from~\eqref{eq: admissible control}, the optimal control process in~\eqref{controlled_dynamics} is characterized as 
\begin{equation}\label{eq: unified_optimal_control}
 u_{\epsilon}(t,x,y,z) = \epsilon \begin{pmatrix}
1 & 0 & -2\alpha y \\
0 & 1 & 2\alpha x
\end{pmatrix} \nabla\log\varphi_\epsilon(t,x,y,z),   
\end{equation}
where $\varphi_\epsilon$ is computed using our explicit heat kernel:
\begin{equation}\label{eq: unified_phi_definition}
\varphi_\epsilon(t,x,y,z) = \int \phi_f(x',y',z') p_{t_f-t,\epsilon}((x,y,z), (x',y',z')) dx'dy'dz',
\end{equation}
for fixed terminal datum $\phi_f$ in~\eqref{eq: SS_BI}. This is the explicit solution of a family of Sub-Riemannian Schr\"odinger Bridges for any $\epsilon>0$.

Since the vector fields  $g_1, g_2$ induces H\"ormander operators on $\mathbb{H}$, from~\cite{BenArous1988,barilari2012small} we have the asympototic formula:
\begin{equation}\label{eq: unified_varadhan}
  -2(t_f - t) \epsilon \log p_{t_f - t,\epsilon}(\vec{q},\vec{q}_1)
  \;\longrightarrow\;
  d_{\mathrm{SR}}^2(\vec{q},\vec{q}_1)
\end{equation}
as $\epsilon\to 0$, where $d_{\mathrm{SR}}$ is the sub-Riemannian distance for the structure defined by the vector fields. For this family of systems, the sub-Riemannian distance has the explicit expression:
\[
d_{\mathrm{SR}}^2(q_1, q_2) = d_{\mathrm{SR}}^2(0, q_1^{-1} q_2)
=
\frac{(2\alpha\theta_c)^2}{\sin^2(2\alpha\theta_c)} (\Delta x^2 + \Delta y^2)
\]
where $\theta_c \in (0,\pi/(2\alpha))$ uniquely solves:
\[
\frac{|\Delta z|}{\Delta x^2 + \Delta y^2} = \frac{1}{2\alpha}\left(\frac{2\alpha\theta_c}{\sin^2(2\alpha\theta_c)} - \cot(2\alpha\theta_c)\right),
\]
with
\[
\Delta x = x_2 - x_1, \quad \Delta y = y_2 - y_1, \quad \Delta z = z_2 - z_1 + 2\alpha(x_1 y_2 - y_1 x_2).
\]

Plugging~\eqref{eq: unified_varadhan} into the representation~\eqref{eq: unified_phi_definition} gives, for small $\epsilon>0$:
\begin{equation}\label{eq: unified_phi_laplace}
  \varphi_\epsilon(t,\vec{q})
  = \int_{\mathbb{H}} \phi_f(\vec{q}_1)
    \exp\left(
      -\frac{d_{\mathrm{SR}}^2(\vec{q},\vec{q}_1)}{2\epsilon(t_f - t)}
      + o\bigl(\tfrac{1}{\epsilon}\bigr)
    \right) d\vec{q}_1,
\end{equation}
and thus by the Laplace principle (see, e.g.,~\cite{FreidlinWentzell2012}), we have that:
\begin{equation}\label{eq: unified_laplace_principle}
  -\epsilon \log \varphi_\epsilon(t,\vec{q})
  \;\longrightarrow\;
   \varphi_0(t,\vec{q}):=\inf_{\vec{q}_1\in\mathbb{H}}
  \left\{
    \Phi_f(\vec{q}_1)
    + \frac{d_{\mathrm{SR}}^2(\vec{q},\vec{q}_1)}{2(t_f - t)}
  \right\},
\end{equation}
where $\Phi_f(\vec{q}_1) := -\lim_{\epsilon\to 0}\epsilon \log \phi_f(\vec{q}_1)$.

Therefore, from~\eqref{eq: unified_optimal_control} and~\eqref{eq: unified_laplace_principle}, we have that:
\[
u_{\epsilon}(t,x,y,z)\;\longrightarrow\;u_0(t,x,y,z):=-\begin{pmatrix}
1 & 0 & -2\alpha y \\
0 & 1 & 2\alpha x
\end{pmatrix} \nabla\varphi_0(t,x,y,z)
\]
is the optimal control for the ensemble of constrained systems~\eqref{eq: unified_dynamics}.}
\end{example}


\bibliographystyle{amsplain}
\bibliography{references}


\end{document}
